\section{Internal Flow (internalSetup)}
\codeword{internalSetup} is a second tool in this repository for closed internal flows, built for Stirling engines (displacer, power piston, heated and cooled walls). It works like \codeword{caseSetup}: an INI-style file in the case directory is merged with the template defaults, the geometry prefixes add their own sections, and the tool writes the OpenFOAM dictionaries, the boundary conditions and the Slurm scripts. Everything specific to external aero (inlet/outlet domain, ride height, cornering, force coefficients) is left out.

The internal templates target ESI OpenFOAM v2412. All the features described here were checked on v2412.

\subsection{Overview}
\begin{itemize}
    \item Entry point: \codeword{internalSetup.py}. Do not run \codeword{caseSetup.py} on an internal case.
    \item Case file: \codeword{internalSetup} (not \codeword{caseSetup}) in the case directory.
    \item Template set: \codeword{setupTemplates/internal}, selected by default (\codeword{-s internal}).
    \item Solver: \codeword{rhoPimpleFoam} (compressible, default) or \codeword{pimpleFoam} (incompressible).
    \item Mesh: \codeword{blockMesh} background box around the geometry + \codeword{snappyHexMesh}. There is no wind tunnel domain, the whole mesh is the gas inside the walls.
    \item Mesh motion: \codeword{displacementInterpolation} along z, driven by a table of body positions against crank angle. The crank angle either follows a prescribed speed, or is solved from the gas forces on the bodies (crank drive, Section \ref{sec:crankDrive}).
    \item No symmetry planes: internal cases are always full models.
\end{itemize}

\subsection{Directory Structure and First Run}
The directory structure is the same as for \codeword{caseSetup}: geometry in \codeword{<job>/02_reference/MSH}, cases in \codeword{<job>/CASES/<case>}. The tool is run from inside the case directory.
\begin{lstlisting}
    cd <job>/CASES/001
    python3 /path/to/caseSetup/internalSetup.py --new
    # edit internalSetup, then
    python3 /path/to/caseSetup/internalSetup.py
\end{lstlisting}
Command line options:
\begin{itemize}
    \item \codeword{--new} - writes a fresh \codeword{internalSetup} file (overwrites an existing one).
    \item \codeword{--modules} - lists the available modules.
    \item \codeword{-d}, \codeword{--controlDict} - rewrites \codeword{system/controlDict} only (e.g. after changing the time step or the run length).
    \item \codeword{-s}, \codeword{--setup} - template set, default \codeword{internal}.
\end{itemize}

As with \codeword{caseSetup}, the first run with a new geometry or a new module fills in the missing keys, writes them to \codeword{internalSetup} and stops with \codeword{updated with default values ... rerun}. Check the new keys and run the tool again. New keys are appended at the end of their section. Sections of parts that are no longer in \codeword{GEOM} are removed from the file.

\textbf{NOTE:} run the setup on the machine where the case will run (e.g. in \codeword{$SCRATCH} on the cluster). The generated scripts hold absolute paths to \codeword{zeroTemplates}, \codeword{meshUtilities} and \codeword{foamLibs} in the repository.

A new \codeword{internalSetup} file looks like this:
\begin{lstlisting}
    [DEFAULTS]
    DEFAULT_MODULES = defaultGlobalControl defaultGlobalSimControl defaultGlobalMaterial defaultGlobalDecomposition defaultMeshMotion
    ADDON_MODULES =

    [TITLES]
    CASENAME = casename
    CASEDESCRP = default description
    JOBCODE = ST

    [GEOMETRY]
    GEOM = WALL-chamber.stl,1,1,3,1.2,low

    [GLOBAL_REFINEMENT]
    BASE_CELL = 0.004
    LOC_IN_MESH = required
    ...
\end{lstlisting}
Remove a module from \codeword{DEFAULT_MODULES} and rerun the tool to expose its settings. While a module is listed there, its template values are used and any section of that name in the file is ignored.

\subsection{Geometry and PIDs}
The \codeword{GEOM} line uses the same columns as \codeword{caseSetup}:
\begin{lstlisting}
    GEOMETRY_FILE_NAME,SCALE,MESH_LEVEL,INFLATION_LAYERS,INFLATION_LAYER_GROWTH_RATE,WALL_MODEL
\end{lstlisting}
Every geometry file needs a prefix that says what the whole body is:
\begin{itemize}
    \item \codeword{WALL} - static walls (chamber, cylinder liner, hot and cold end faces).
    \item \codeword{MOVE} - a rigid moving body (displacer). A body whose faces are part of another file (e.g. a piston crown in the chamber file) has a \codeword{MOVE} section without a file of its own, see \codeword{FROM} below.
\end{itemize}
\codeword{REGN} (regenerator) and \codeword{SMP} (sample surfaces) are reserved prefixes and are not available yet.

Inside each file the faces are separated into PIDs: STL \codeword{solid} names or OBJ groups. Every PID must have a name. \codeword{snappyHexMesh} keeps every PID as its own patch, named \codeword{<fileName>_<pid>} (e.g. \codeword{WALL-chamber_hot}). In the setup file the PIDs are selected by keyword: a keyword matches every PID that contains it, ignoring case. The tool reads the PIDs directly from the geometry files, checks the keywords before anything is written and prints a table of every patch with its group, U and T condition.

\textbf{NOTE:} the mesh is the gas side of the walls, chosen by \codeword{LOC_IN_MESH} in \codeword{[GLOBAL_REFINEMENT]}. This point is required and must be inside the gas at the meshing (start) position of the bodies.

\subsection{Wall Boundary Conditions}
Each \codeword{WALL} and \codeword{MOVE} geometry gets a section with the same name as the file (without the extension):
\begin{lstlisting}
    [WALL-chamber]
    U_DEFAULT = noSlip
    T_DEFAULT = adiabatic
    BC_GROUPS = hotEnd coldEnd liner
    hotEnd = PIDS hot T fixed 600
    coldEnd = PIDS cold T fixed 300
    liner = PIDS liner T convective 15 300
    ANCHOR_PIDS = hot cold
    GAP_PIDS = side
\end{lstlisting}
\begin{itemize}
    \item \codeword{U_DEFAULT} [\codeword{str}] - velocity condition for PIDs that are in no group (\codeword{noSlip}, \codeword{slip}).
    \item \codeword{T_DEFAULT} [\codeword{str}] - temperature condition for PIDs that are in no group, options as in the group lines below.
    \item \codeword{BC_GROUPS} [\codeword{str}] - names of the groups. Each name needs its own line:
    \begin{lstlisting}
    <name> = PIDS <keyword> [<keyword> ...] [U noSlip|slip] [T <mode> <values>]
    \end{lstlisting}
    Temperature modes:
    \begin{itemize}
        \item \codeword{adiabatic} - zero gradient.
        \item \codeword{fixed <T>} - fixed temperature in K.
        \item \codeword{convective <h> <Ta>} - heat transfer coefficient (W/m$^2$K) to an ambient temperature (K).
        \item \codeword{flux <q>} - fixed heat flux in W/m$^2$.
    \end{itemize}
    A PID that matches two groups and a keyword that matches no PID are errors. PIDs in no group form the group \codeword{<file>_default}.
    \item \codeword{ANCHOR_PIDS} [\codeword{str}] - flat faces that hold the mesh still for the motion solver, normally the hot and cold end faces (\codeword{WALL} only). Add any fixed step or port end that lies between an anchor and a moving body, otherwise its faces move with the mesh. The setup warns about such faces.
    \item \codeword{GAP_PIDS} [\codeword{str}] - wall side of the displacer side gap. Used for the gap refinement, see Section \ref{sec:motion}.
\end{itemize}
Faces that move with a body always get \codeword{movingWallVelocity}, whatever U is set to. In an incompressible run (\codeword{pimpleFoam}) there is no temperature field and the T settings are ignored (the setup prints a note).

\subsection{Run Control}
\begin{lstlisting}
    [GLOBAL_CONTROL]
    STARTFROM = latestTime
    FREQ = 10
    N_CYCLES = 3
    MAX_CO = 1
    MAX_DT = 1e-4
    INIT_DT = 1e-6
    WRITE_PER_CYCLE = 20
\end{lstlisting}
\begin{itemize}
    \item \codeword{FREQ} [\codeword{float,Hz}] - engine speed. With the crank drive it is the start speed.
    \item \codeword{N_CYCLES} [\codeword{float}] - number of cycles to run, \codeword{endTime = N_CYCLES/FREQ}.
    \item \codeword{WRITE_PER_CYCLE} [\codeword{int}] - number of written times per cycle.
    \item \codeword{MAX_CO} [\codeword{float}] - Courant number limit for the adjustable time step. Set it to \codeword{0} or leave it blank (\codeword{MAX_CO =}) for a fixed time step equal to \codeword{MAX_DT} (\codeword{adjustTimeStep no}). With a fixed step the setup prints the number of steps per cycle, and warns if the write interval is not a whole number of steps (the writes then happen at the first step past each interval).
    \item \codeword{MAX_DT} [\codeword{float,s}] - largest time step (adjustable), or the time step (fixed).
    \item \codeword{INIT_DT} [\codeword{float,s}] - first time step of an adjustable run, not used with a fixed step.
    \item \codeword{STARTFROM} [\codeword{str}] - \codeword{latestTime} so a resubmitted run continues.
\end{itemize}
After changing these values, \codeword{internalSetup.py -d} rewrites the controlDict only.

\subsection{Solver and Gas}
\begin{lstlisting}
    [GLOBAL_SIM_CONTROL]
    SOLVER = rhoPimpleFoam
    TURB_MODEL = laminar
    DISSIPATION = no
    N_OUTER_CORR = 1
    N_CORR = 2
    N_NONORTH_CORR = 1
    P_MAX_ITER = 1000
    RELAX_P = auto
    RELAX_U = auto
    RELAX_E = auto

    [GLOBAL_MATERIAL]
    GAS = air
    CHARGE_PRESSURE = 1e5
    INIT_T = 300
    INIT_T_FIELD = uniform
    TRANSPORT = const
\end{lstlisting}
\begin{itemize}
    \item \codeword{SOLVER} - \codeword{rhoPimpleFoam}: perfect gas, absolute pressure in Pa, energy equation, wall heat transfer. \codeword{pimpleFoam}: incompressible, kinematic pressure (p/$\rho$) referenced to 0 at \codeword{LOC_IN_MESH}, no temperature. The incompressible option is only useful for looking at the flow through the displacer: the gas is not compressed, so the pressure level, the piston work and the heat transfer have no meaning.
    \item \codeword{TURB_MODEL} - \codeword{laminar} (default), \codeword{kOmegaSST}, \codeword{kEpsilon} or \codeword{kOmega}. The wall model column of \codeword{GEOM} (\codeword{low}/\codeword{high}) selects low or high Reynolds wall treatment.
    \item \codeword{DISSIPATION} - \codeword{no} (default) or \codeword{yes}. \codeword{rhoPimpleFoam} solves the energy in total-energy form: $\partial(\rho e)/\partial t + \nabla\cdot(\rho \mathbf{U} e) + \partial(\rho K)/\partial t + \nabla\cdot(\rho \mathbf{U} K) + \nabla\cdot(p\mathbf{U}) = -\nabla\cdot\mathbf{q}$, with $K = |\mathbf{U}|^2/2$ and without the viscous work $\nabla\cdot(\boldsymbol{\tau}\cdot\mathbf{U})$ and the work of momentum sources. Subtracting the kinetic energy equation (momentum $\cdot\,\mathbf{U}$) leaves $\rho\,De/Dt = -p\nabla\cdot\mathbf{U} - \mathbf{U}\cdot\nabla\cdot\boldsymbol{\tau} - \mathbf{U}\cdot\mathbf{f} - \nabla\cdot\mathbf{q}$: the drag work $-\mathbf{U}\cdot\mathbf{f}$ of a porous zone is in $e$, and $-\mathbf{U}\cdot\nabla\cdot\boldsymbol{\tau} = \boldsymbol{\tau}:\nabla\mathbf{U} - \nabla\cdot(\boldsymbol{\tau}\cdot\mathbf{U})$ is the viscous dissipation up to a divergence, which over the whole gas is only the viscous work of the walls. The heat made by friction and drag is therefore already in the gas energy; only its placement near walls differs. \codeword{yes} adds the ESI \codeword{viscousDissipation} fvOption and the drag work of every porous zone (\codeword{porDragHeat_<zone>}) on top, which counts that heat twice (seen in a closed-engine test as a gas energy rise above the heat and work put in). It is kept for testing only.
    \item \codeword{N_OUTER_CORR} - PIMPLE outer correctors. \codeword{1} is PISO: stable only around a Courant number of 1 on the moving mesh, so keep \codeword{MAX_CO} near 1. For larger time steps (\codeword{MAX_CO} well above 1) use 3 or more outer correctors, without under-relaxation (below).
    \item \codeword{P_MAX_ITER} - most iterations of the pressure solver per solve (default 1000, the OpenFOAM default). In a near-static gas, such as a very slow piston, the pressure solver cannot get below about $10^{-4}$ (normalised residual, a round-off floor) and every solve then runs all the iterations, many times slower without changing the result. Lower it to about 50 for such cases; the mass check (\codeword{gasMass}) shows whether the run still conserves mass.
    \item \codeword{RELAX_P}, \codeword{RELAX_U}, \codeword{RELAX_E} - under-relaxation of the pressure field and of the U (and k, omega) and e equations. \codeword{auto} (default): 0.3, 0.7, 0.7 for the steady solvers, none (1) for the transient ones. The last outer iteration of each time step is never relaxed. Do not relax a transient compressible run: with the pressure relaxed, the density $\psi p$ lags the continuity equation in the relaxed outer iterations, and the energy equation of the last iteration still uses that density, so the gas energy is wrong by the lag times $\rho e$ in every step. In a closed-cylinder compression test with 3 outer correctors and p relaxed by 0.3, half of the compression heating was lost (the gas ended 15\,K too cold after a 23\,\% volume change); without relaxation the energy balance closed within 0.2\,\% at the same cost, and with relaxation it took 20 outer correctors to get within 1\,\%. The setup warns when a transient compressible run is relaxed.
    \item \codeword{GAS} - \codeword{air}, \codeword{nitrogen}, \codeword{helium}, \codeword{hydrogen} or \codeword{custom}. Each preset sets the molecular weight, Cp, viscosity and Prandtl number. Any of \codeword{MOL_WEIGHT}, \codeword{CP}, \codeword{MU}, \codeword{PR} that is not \codeword{default} overrides the preset, \codeword{custom} needs all four.
    \item \codeword{CHARGE_PRESSURE} [\codeword{float,Pa}], \codeword{INIT_T} [\codeword{float,K}] - initial gas state. The gas mass follows from these and the meshed volume.
    \item \codeword{INIT_T_FIELD} - \codeword{uniform} (default): the gas starts at \codeword{INIT_T}. \codeword{matrix} (the only porous zone with two \codeword{SOLID_T} values) or a \codeword{POR} zone name: warm start, the gas temperature follows the matrix start profile along $z$ (first \codeword{SOLID_T} value at the $-z$ end of the zone, second at the $+z$ end, linear in between) and stays at the end values beyond the zone, so the hot and cold spaces start at their wall temperatures. The pressure is raised to a uniform value that keeps the gas mass of the charge state (charged at \codeword{CHARGE_PRESSURE} and \codeword{INIT_T}, then heated): $p = \sum p_c V_i/T_c \,/\, \sum V_i/T_i$. The run script does it with \codeword{remeshMassFix -warmStart} after the zero directories are made, only while the case has no later time, and prints the new mean temperature and pressure in \codeword{log.warmStart}. A warm start avoids the first cycles in which a cold gas and a warm matrix only even each other out.
    \item \codeword{TRANSPORT} - \codeword{const} viscosity or \codeword{sutherland} (coefficients from the preset, or \codeword{SUTHERLAND_AS} / \codeword{SUTHERLAND_TS}).
    \item \codeword{NU} - kinematic viscosity for \codeword{pimpleFoam}, \codeword{default} = $\mu/\rho$ of the gas at the charge state.
\end{itemize}

\subsection{Mesh Motion}\label{sec:motion}
A case with one or more \codeword{MOVE} sections gets a moving mesh. Static cases (only \codeword{WALL}) skip all of this.

\subsubsection{How the Mesh Moves}
The motion solver is \codeword{displacementInterpolation} along z. Flat face zones normal to z are made from the anchor faces (fixed) and from the moving faces of each body (moving with that body). The displacement of a mesh point depends only on its z coordinate: points between two zones of the same body move rigidly with it (the displacer, its holes and the side gap), and points between zones of different bodies, or between a body and an anchor, are stretched or squashed linearly (the end spaces and the piston bore). Points beyond the outermost zone interpolate towards zero at the mesh bound, which is fine when the end faces are the mesh bounds in z.

This works for a displacer and a power piston in line along z (gamma or beta type engines), as long as every zone is its own flat z band and the bands never cross. Side-by-side cylinders overlapping in z with different motions (alpha type) are not supported.

\subsubsection{Motion Table}
Both bodies are driven by one CSV file with the shared crank angle, so the phase between the bodies is fixed by the file:
\begin{lstlisting}
    theta,displacer,piston
    0,0.0,0.0
    1,0.0000012,-0.0000017
    ...
    359,...
\end{lstlisting}
\begin{itemize}
    \item \codeword{theta} in degrees, starting at 0 and ending at or below 360, any spacing. Lines starting with \codeword{#} are comments.
    \item One column per body: axial position in metres, any datum (only differences are used).
    \item The profile must be periodic: a value given at 360 must equal the value at 0.
\end{itemize}

\subsubsection{MOVE Sections}
\begin{lstlisting}
    [MOVE-displacer]
    FROM = self
    MOTION_FILE = ../../02_reference/engineMotion.csv
    MOTION_COLUMN = displacer
    REF_ANGLE = 0
    TDC_CLEARANCE = 0.002
    ZONE_PIDS = top bottom
    GAP_PIDS = side
    GAP_CELLS = 6

    [MOVE-piston]
    FROM = WALL-chamber
    MOTION_FILE = ../../02_reference/engineMotion.csv
    MOTION_COLUMN = piston
    REF_ANGLE = 0
    TDC_CLEARANCE = 0.002
    ZONE_PIDS = crown
\end{lstlisting}
\begin{itemize}
    \item \codeword{FROM} [\codeword{str}] - \codeword{self}: the body is its own \codeword{MOVE-<name>} file and the whole file moves. \codeword{<geometry>}: the body is made of PIDs in another file (e.g. the piston crown, a PID of \codeword{WALL-chamber}). Only its \codeword{ZONE_PIDS} faces move. A \codeword{FROM} section stays in the file as long as that geometry is in \codeword{GEOM}. The crown must be a separate, flat shell normal to z, slightly larger than the bore so that it seals.
    \item A PID named \codeword{MOVE-<name>} in a WALL file (e.g. \codeword{MOVE-piston}) creates the section \codeword{[MOVE-<name>]} automatically: \codeword{FROM} = that file, \codeword{ZONE_PIDS} = the PID, \codeword{MOTION_COLUMN} = \codeword{<name>}, the other keys from the defaults, and the setup stops so it can be checked. Each such PID is its own body. When the PID is no longer in the file, the section is removed. A \codeword{ZONE_PIDS} keyword that is a whole PID name only takes that PID.
    \item \codeword{MOTION_FILE} [\codeword{path}] - motion CSV, required. Absolute, with \codeword{~} and \codeword{$VARIABLES} expanded, or relative to the case directory.
    \item \codeword{MOTION_COLUMN} [\codeword{str}] - column of this body in the CSV.
    \item \codeword{REF_ANGLE} [\codeword{float,deg}] - crank angle at which the geometry was exported. Before meshing, the tool moves the body from this angle to the start angle.
    \item \codeword{TDC_CLEARANCE} [\codeword{float,m}] - smallest allowed gap to a neighbouring anchor (the hot or cold face) over the cycle. A gap of zero cannot be meshed. How a profile that comes closer is changed is set by \codeword{CLEARANCE_MODE}.
    \item \codeword{CLEARANCE_MODE} - \codeword{soft} (default): only the end of the stroke that comes too close is eased back, with the smooth floor $z' = z_{end} + ((z - z_{end})^4 + \delta^4)^{1/4}$; the rest of the CSV profile is kept (for a 2 mm clearance the change is below 0.05 mm from 5 mm away from that end), dwells stay flat, both ends are handled if both are too close, and the far end stays exactly where the CSV puts it (so a piston crown stays at its liner end). The peak acceleration near that end rises (about doubles for a 2 mm clearance on a 152 mm stroke). \codeword{scale}: the whole stroke is scaled down between the limits, the shape is kept exactly but all velocities drop. \codeword{shift}: the whole path and the geometry are offset by the shortfall (stroke kept, both ends move). The setup prints the stroke before and after, the largest change from the CSV and the added dead volume.
    \item Sealing: a body made of PIDs in a WALL file (\codeword{FROM}, the piston crown) must, at the meshing angle, lie inside its liner along z and cover the bore (the liner is the non-flat wall in the same file with about the crown's footprint), else the setup stops before meshing and says how much liner is missing or that the crown is smaller than the bore. After meshing the liner boundary stretches with the mesh, so only the meshing angle matters.
    \item \codeword{ZONE_PIDS} [\codeword{str}] - flat faces that drive the motion: displacer \codeword{top bottom}, piston \codeword{crown}.
    \item \codeword{GAP_PIDS} [\codeword{str}], \codeword{GAP_CELLS} [\codeword{int}] - body side of the side gap and the minimum number of cells across it (default 6, 4 for quick tests). The tool measures the gap from the STLs and refines a band around it with a \codeword{snappyHexMesh} distance refinement. Gap faces get no inflation layers. Needs \codeword{GAP_PIDS} on the \codeword{WALL} side as well.
    \item \codeword{U_DEFAULT}, \codeword{T_DEFAULT}, \codeword{BC_GROUPS} - as for \codeword{WALL}.
\end{itemize}

\subsubsection{MESH\_MOTION Settings}
\begin{lstlisting}
    [MESH_MOTION]
    METHOD = displacementInterpolation
    AXIS = z
    START_ANGLE = auto
    STEP_DEG = 1
    MAX_DEFORMATION = 15
    REMESH = no
    REMESH_DEFORMATION = 3
    GAP_CELLS_Z = 6
    REMESH_DT = 1e-6
    ZONE_FLAT_TOL = 1e-4
    MOTION_FRAMES = 24
    DRIVE = prescribed
\end{lstlisting}
\begin{itemize}
    \item \codeword{START_ANGLE} [\codeword{float,deg}] - crank angle of the mesh and of the start of the run. \codeword{auto} picks the angle that gives the smallest worst squash/stretch of the gaps over the cycle.
    \item \codeword{STEP_DEG} [\codeword{float,deg}] - resampling of the profile ($\le$ 1).
    \item \codeword{MAX_DEFORMATION} [\codeword{float}] - largest allowed ratio between a gap at the meshing angle and the same gap anywhere in the cycle (both ways). Larger: increase \codeword{TDC_CLEARANCE}, or remesh. Not used with remeshing.
    \item \codeword{REMESH} - \codeword{no} (one mesh for the whole run) or \codeword{auto} (a new mesh whenever a gap has been squashed or stretched \codeword{REMESH_DEFORMATION} times, Section \ref{sec:remesh}).
    \item \codeword{REMESH_DEFORMATION} [\codeword{float}] - squash/stretch allowed per mesh with \codeword{REMESH = auto} (above 1).
    \item \codeword{GAP_CELLS_Z} [\codeword{float}] - smallest number of cells across a gap between two zones (end spaces) at each meshing angle. Thinner gaps get a band refined in z only (flat cells, Section \ref{sec:remesh}). 0 = off.
    \item \codeword{REMESH_DT} [\codeword{float,s}] - time step of the first step after each map (default $10^{-6}$), then the adjustable time step ramps up again.
    \item \codeword{ZONE_FLAT_TOL} [\codeword{float,m}] - how flat a zone face must be.
    \item \codeword{DRIVE} - \codeword{prescribed} (the crank turns at \codeword{FREQ}) or \codeword{crank} (Section \ref{sec:crankDrive}).
\end{itemize}

\subsubsection{Static Walls in a Moving Mesh}
The points of static walls between the motion zones (side wall, liner) slide along the wall with the mesh. Snapped faces are not exactly parallel to the motion, so they sweep a little volume each step. With a plain no-slip wall that volume leaks through the wall (of the order of 1\,\% of the gas mass per cycle). Static no-slip walls of a moving-mesh case therefore get \codeword{slidingWallVelocity} (\codeword{foamLibs/stirlingMotion}): the velocity normal to the face matches the mesh motion and the tangential velocity is zero, so no gas crosses the wall and the wall does not drag the gas. Faces that do not move (the anchors) get exactly zero. The case loads the library in \codeword{controlDict} and the run scripts compile it. Walls set to \codeword{slip} are not covered yet.

\subsubsection{Setup Checks and Report}
Before writing the case, the setup checks the zones over every angle of the cycle (not only the start): each zone is flat, the z order stays the same with a gap above zero, the clearances hold, and the deformation stays below \codeword{MAX_DEFORMATION}. It prints the zones in z order, the start angle, the worst squash/stretch, and for each body the stroke, peak velocity and peak acceleration. A sharp step in the profile is reported, since it gives pressure spikes. Use the peak velocity to choose \codeword{MAX_DT}.

\subsubsection{Mesh Motion Dry Run}
Check the moving mesh before the flow run:
\begin{lstlisting}
    ./runScript --meshOnly
    ./motionCheckScript
    ./postScript --mesh
\end{lstlisting}
or on the cluster \codeword{./submitScript --check}, which chains all three. The motion check runs \codeword{moveDynamicMesh} over one cycle and \codeword{checkMesh} on the moved meshes every 45 degrees, and prints the quality summary. It stops only if \codeword{moveDynamicMesh} fails. Mesh quality failures are printed for the user to judge.

Before removing the moved meshes, it exports them to \codeword{EnSightMotion}, \codeword{MOTION_FRAMES} times per cycle (\codeword{[MESH_MOTION]}, a multiple of 8, default 24 = every 15 degrees; the motion itself still steps 5 degrees or less). \codeword{./postScript --mesh} then draws the slices of \codeword{pvPostInternalSetup} (\codeword{[SLICE]}, including \codeword{FOLLOW}) and the body surfaces (\codeword{[SURFACE] PATCHES} and \codeword{VIEWS}) as grey surface-with-edges at every exported crank angle, in \codeword{postProcessing/meshImages}, and makes one video per panel in \codeword{postProcessing/meshMovies}. Watch the end spaces squash and stretch, the displacer band move without deforming, and the cells next to the anchors.

\textbf{NOTE:} \codeword{EnSightMotion} holds the whole mesh for every frame, so its size grows with the cell count times \codeword{MOTION_FRAMES}. Delete it once the videos have been checked.

\subsection{Crank Drive}\label{sec:crankDrive}
With \codeword{DRIVE = crank} the gas forces turn the crank instead of a prescribed speed. The bodies stay tied to the crank angle $\theta$ through the motion CSV; only $\theta(t)$ is solved. A custom motion solver, \codeword{crankDisplacementInterpolation} in \codeword{foamLibs/stirlingMotion}, does this. The run scripts compile it with \codeword{wmake libso} into \codeword{$FOAM_USER_LIBBIN} (this does nothing when it is up to date), so a compiler must be available where the job runs.

Every time step:
\begin{equation}
F_{body} = \sum (p - p_{ref})\, \mathbf{S}_f \cdot \mathbf{e}_z, \qquad T_{body} = F_{body}\,\frac{dz}{d\theta}
\end{equation}
\begin{equation}
I\,\frac{d\omega}{dt} = \sum T_{body} - T_{load} - T_{friction}
\end{equation}
The forces are pressure only, from the previous time step. Speed-proportional terms are treated implicitly. The full equations, the load models and the summary values are in Section~\ref{sec:postMath}.
\begin{lstlisting}
    [MESH_MOTION]
    DRIVE = crank
    INERTIA = 0.05
    LOAD = generator
    LOAD_TORQUE = 0
    LOAD_DAMPING = 0
    GEN_K = 1
    GEN_R_LOAD = 10
    GEN_R_INT = 0
    FRICTION_TORQUE = 0
    FRICTION_COEFF = 0
    P_BUFFER = charge
\end{lstlisting}
\begin{itemize}
    \item \codeword{INERTIA} [\codeword{float,kg m2}] - flywheel, crank and everything turning with it.
    \item \codeword{LOAD} - what the crank drives:
    \begin{itemize}
        \item \codeword{speed} - fixed speed at \codeword{FREQ} (same motion as \codeword{prescribed}, the gas torque is logged).
        \item \codeword{torque} - constant load torque \codeword{LOAD_TORQUE} (N m).
        \item \codeword{damper} - torque \codeword{LOAD_DAMPING} $\cdot\,\omega$ (N m s).
        \item \codeword{generator} - permanent magnet generator into a resistor: $I = k\omega/(R_{load} + R_{int})$, torque $kI$, electrical power $I^2 R_{load}$. \codeword{GEN_K} (V s/rad), \codeword{GEN_R_LOAD} and \codeword{GEN_R_INT} (ohm).
    \end{itemize}
    \item \codeword{FRICTION_TORQUE} [\codeword{N m}] (Coulomb) and \codeword{FRICTION_COEFF} [\codeword{N m s}] (viscous).
    \item \codeword{P_BUFFER} - pressure behind the bodies, \codeword{charge} = \codeword{CHARGE_PRESSURE}, or a value in Pa. The bounce space is not modelled.
\end{itemize}
The setup prints the load torque and power at the start speed. The crank angle and speed are stored in \codeword{<time>/uniform/crankState} at each write, so a restart continues from them. The log \codeword{postProcessing/crank/<start>/crank.dat} has the columns: time, crank angle (accumulated, deg), $\omega$, gas torque, load torque, electrical power, and the force and torque per body.

\textbf{NOTE:} \codeword{endTime} is still \codeword{N_CYCLES/FREQ}. When the speed changes, the run covers about \codeword{N_CYCLES} cycles, not exactly. With \codeword{pimpleFoam} the crank drive only makes sense with \codeword{LOAD = speed} (the setup warns).

\subsection{Remeshing}\label{sec:remesh}
With \codeword{REMESH = auto} the cycle is split into segments. Each segment has its own mesh, made at the crank angle where the segment starts, and the mesh only squashes or stretches until a gap has changed \codeword{REMESH_DEFORMATION} times (one way). The setup walks the cycle from the start angle in \codeword{STEP_DEG} steps and starts a new segment at the last angle within the limit, so the segments are the same every cycle and each mesh is made once. It prints the segments and their worst squash/stretch, and checks the crown seal and \codeword{LOC_IN_MESH} at every meshing angle (\codeword{LOC_IN_MESH} must be in the gas at all of them).
\begin{itemize}
    \item The case mesh is the first segment (\codeword{seg01}). The others are full cases in \codeword{meshes/seg02 ...} with the geometry moved to their start angle. \codeword{system/remeshSegments} lists the mesh names with their start and end angles.
    \item \codeword{runScript} makes all meshes in the same job (snappy, topoSet, zone reordering and checkMesh for each), then loops: solve until the crank reaches the end of the segment, map the fields onto the next mesh with \codeword{mapFieldsPar}, restart. The map is volume weighted (conservative); the few new cells along the walls that the old mesh does not overlap at all (snapped walls never match exactly) would be left at zero, so \codeword{remeshMassFix -combine} fills them from their neighbours on the new mesh; the wall values are mapped from the nearest old face. A separate nearest-cell map is not used: it crashes in v2412 on large parallel cases. The crank motion solver stops the run itself (\codeword{stopAngle} in each \codeword{constant/dynamicMeshDict.seg*}); its last step lands exactly on the segment end. Prescribed motion runs as the crank drive with \codeword{LOAD = speed}.
    \item The new mesh is written into the time directory of the map (\codeword{polyMesh} with \codeword{points0} and a \codeword{remeshSegment} marker), so the run stays one case and the post-processing works as before.
    \item Mapping changes the gas mass slightly (of the order of $10^{-4}$ per map). For \codeword{rhoPimpleFoam} the utility \codeword{remeshMassFix} (\codeword{foamLibs/remeshMassFix}, compiled by the run script) scales p and $\rho$ back to the mass before the map, T is kept.
    \item Surface fluxes cannot be mapped, so the restart rebuilds them from U. The first step after a map uses \codeword{REMESH_DT} ($10^{-6}$ s) and the adjustable time step ramps up again (needs \codeword{MAX_CO} above 0). With $10^{-5}$ s one map of an engine cycle diverged in its first step.
    \item A mesh made near a stroke end has a thin end space (a few mm between the displacer and the end face). Every gap between two zones that would get fewer than \codeword{GAP_CELLS_Z} cells across at the meshing angle gets a band over the gap that snappy refines in z only (directional refinement, \codeword{levelIncrement}), so the cells there are flat, as in a deformed mesh, instead of one or two cubes. \codeword{LOC_IN_MESH} of each segment mesh is moved to the middle of the gap it lies in (snappy picks the mesh region before the z refinement). The setup prints the bands per meshing angle; each segment has its own \codeword{meshes/segNN/snappyHexMeshDict}.
    \item The new mesh files are hard links to the segment mesh (no copy), so a map takes about a minute.
    \item The motion dry run (\codeword{motionCheckScript}) only covers the first segment.
\end{itemize}
Not supported yet with remeshing: \codeword{FAN} baffles.

\subsection{Decomposition and Cluster Scripts}
\begin{lstlisting}
    [GLOBAL_COMPUTE_SETUP]
    NUM_PROCS = 192
    LOCAL_PROCS = 4
    DECOMP_METHOD = scotch
    DECOMP = auto
    NODELIST =
    NODES = 1
    EXCLUSIVE = True
    CPUS_PER_TASK = 1
    MEM_PER_CPU = 4G
    JOB_NAME = auto
    RUN_TIME = 0-12:00
    MOTION_TIME = 0-02:00
    EXPORT = ensight
    POST_CPUS = 8
    POST_TIME = 0-03:00

    [CLUSTER_ENVIRONMENT]
    ACCOUNT = def-xxxx
    MODULES = StdEnv/2023 gcc/12.3 openmpi/4.1.5 openfoam/v2412
    MPI_RUN = srun
    PARTITION =
    POST_MODULES = StdEnv/2023 python/3.11 scipy-stack paraview/5.13.1 ffmpeg
    POST_MEM = 32G
    POST_ENV =
    POST_CONFIG =
    PVBATCH = pvbatch
    GSHEET = True
    COPY_BACK = True
\end{lstlisting}
\codeword{GLOBAL_COMPUTE_SETUP} holds the settings of a run, \codeword{CLUSTER_ENVIRONMENT} those of the cluster it runs on. \codeword{CLUSTER_ENVIRONMENT} is in \codeword{DEFAULT_MODULES}, so it does not appear in \codeword{internalSetup}: take \codeword{defaultClusterEnvironment} out of \codeword{DEFAULT_MODULES} to change it. A case file from before this split keeps its cluster keys in \codeword{GLOBAL_COMPUTE_SETUP}; they are still used (the setup prints a note).
\begin{itemize}
    \item \codeword{NUM_PROCS} - Slurm tasks (192 = one full node on a 192-core cluster). \codeword{LOCAL_PROCS} - cores for a local run.
    \item \codeword{DECOMP_METHOD} - \codeword{scotch} (default), or \codeword{hierarchical} / \codeword{simple} for clusters without scotch.
    \item \codeword{DECOMP} - split for \codeword{hierarchical}/\codeword{simple}. \codeword{auto}: the scripts choose \codeword{n (nx ny nz)} for the task count of each run, with subdomains of the background mesh box as close to cubes as possible. \codeword{nx ny nz}: that split is used, and \codeword{nx}$\cdot$\codeword{ny}$\cdot$\codeword{nz} must equal \codeword{NUM_PROCS} (and \codeword{LOCAL_PROCS} for a local run). Ignored for other methods.
    \item \codeword{MPI_RUN} - how the parallel tools are started inside a Slurm job: \codeword{srun} (default), or e.g. \codeword{/usr/bin/mpirun} on clusters where Open MPI was not built to start under srun (all ranks then abort in \codeword{MPI_Init}). Outside Slurm the scripts always use \codeword{mpirun -np LOCAL_PROCS}.
    \item \codeword{NODES}, \codeword{NODELIST} (a fixed node, \codeword{--nodelist}), \codeword{JOB_NAME} (\codeword{auto} = case directory name), \codeword{RUN_TIME}, \codeword{MOTION_TIME} - Slurm header of the run.
    \item \codeword{EXCLUSIVE} - \codeword{True} (default): the run and motion-check jobs take whole nodes with all their memory (\codeword{--exclusive}, \codeword{--mem=0}). \codeword{False}: the jobs share the node, for small meshes on a few cores (aim for about 5000 to 10000 cells per core; far more cores than that makes a run slower, not faster). A shared job asks for \codeword{MEM_PER_CPU} per core (e.g. \codeword{4G}; blank: the cluster's default, which can be small).
    \item \codeword{CPUS_PER_TASK} - cores given to each MPI rank (default 1). OpenFOAM runs one thread per rank and is limited by memory bandwidth, so on large nodes fewer ranks spread evenly over the node can be nearly as fast as one rank per core. With \codeword{NUM_PROCS = 48} and \codeword{CPUS_PER_TASK = 4} on a 192-core node, the ranks sit two per 8-core group (shared L3 cache) instead of filling the first 48 cores. Works with \codeword{MPI_RUN = srun} (the scripts pass it to srun with \codeword{SRUN_CPUS_PER_TASK}); \codeword{MEM_PER_CPU} counts every core, so the memory per rank is \codeword{CPUS_PER_TASK} times it.
    \item \codeword{EXPORT} - \codeword{ensight} (\codeword{foamToEnsight} at the end of the run) or \codeword{none}. The images and videos need the EnSight export.
    \item \codeword{POST_CPUS}, \codeword{POST_TIME} - the post-processing job (one task with \codeword{POST_CPUS} cores).
    \item \codeword{ACCOUNT}, \codeword{MODULES}, \codeword{PARTITION} - Slurm account, environment modules and partition. Blank drops the line (clusters without accounts or modules, where OpenFOAM comes from the shell that submits).
    \item \codeword{POST_MODULES}, \codeword{POST_MEM} - modules and memory of the post job. The modules must provide \codeword{python3} with numpy, \codeword{pvbatch} and \codeword{ffmpeg}. \codeword{POST_MEM} blank: no memory request (nodes without a memory size in the Slurm configuration reject any \codeword{--mem}).
    \item \codeword{POST_ENV} - a Python virtual environment folder that \codeword{postScript} activates for the gsheet and Drive steps only, e.g. \codeword{~/gdrive_env} with \codeword{gspread} and \codeword{google-auth} installed (blank: none). \codeword{POST_CONFIG} - the folder with the gsheet and Drive settings (blank: \codeword{~/.config/caseSetup}); on clusters whose compute nodes do not share the home directory it has to be on a shared disk.
    \item \codeword{PVBATCH} - the pvbatch command (a full path for a local ParaView install; without a display it must be an offscreen build, e.g. OSMesa or EGL).
    \item \codeword{GSHEET} - \codeword{True}: \codeword{postScript} pushes the summary to the database sheet (needs \codeword{gspread} and the sheet settings, see Section \ref{sec:internalPost}). \codeword{COPY_BACK} - \codeword{True}: \codeword{postScript} copies the run data to Google Drive.
\end{itemize}
The scripts set \codeword{numberOfSubdomains} (and the hierarchical split) every time they run, so the same case can be run on the cluster and locally.

Generated scripts (each runs as a Slurm job with \codeword{sbatch}, or locally with \codeword{./script} using \codeword{mpirun}):
\begin{itemize}
    \item \codeword{runScript} - mesh + solve + export in one job: \codeword{blockMesh}, \codeword{surfaceFeatureExtract}, \codeword{decomposePar}, \codeword{snappyHexMesh}, \codeword{topoSet} (motion zones), face zone reorder, \codeword{checkMesh}, \codeword{createZeroDirectory}, the solver and \codeword{foamToEnsight}. The mesh is kept (\codeword{.meshDone}): running it again skips the meshing and continues the solve from the latest time. Remove \codeword{processor*} and \codeword{.meshDone} to remesh. The setup writes \codeword{system/meshSignature} (a checksum of the meshing dicts, the motion zones in \codeword{topoSetDict} and the geometry as meshed). If rerunning the setup changes it, \codeword{runScript} and \codeword{motionCheckScript} stop and ask for a remesh, instead of running on a mesh that no longer matches the setup (for example a new MOVE body whose motion zone was never made). \codeword{--meshOnly} stops after the mesh. Earlier solver logs are kept as \codeword{log.<solver>.N}.
    \item \codeword{motionCheckScript} - the mesh motion dry run (moving cases only).
    \item \codeword{postScript} - post-processing (Section \ref{sec:internalPost}): summary, gsheet row, images and videos. A failed step is reported and the next step still runs; the job then ends with exit code 1.
    \item \codeword{submitScript} - submits \codeword{runScript} and then \codeword{postScript} (started when \codeword{runScript} has finished without error). \codeword{--noPost}: \codeword{runScript} only. \codeword{--post}: \codeword{postScript} only, e.g. after changing the image settings. \codeword{--check}: \codeword{runScript --meshOnly} followed by \codeword{motionCheckScript}.
\end{itemize}

\textbf{NOTE:} v2412 has a bug in \codeword{displacementInterpolation}: if the motion zones are not the first face zones in the mesh, every zone silently gets the wrong displacement. The scripts run \codeword{meshUtilities/reorderFaceZones.py} after \codeword{topoSet} to put them first. The crank drive solver matches zones by name and does not have this problem.

\subsection{Through-Flow: Inlets, Outlets and Steady Solvers}\label{sec:throughFlow}
For pipes, ducts and exhaust stacks the gas enters and leaves the domain. Inlets and outlets are BC groups on PIDs, in the same \codeword{WALL} section as the wall groups:
\begin{lstlisting}
    [WALL-stack]
    U_DEFAULT = noSlip
    T_DEFAULT = adiabatic
    BC_GROUPS = stackIn stackOut stackWall
    stackIn = PIDS inlet U massFlow 0.5 T fixed 600
    stackOut = PIDS outlet P fixed 1e5
    stackWall = PIDS wall T convective 10 300
\end{lstlisting}
\begin{itemize}
    \item \codeword{U massFlow <kg/s>} - mass flow inlet (into the domain).
    \item \codeword{U velocity <m/s>} - velocity normal to the patch, into the domain. \codeword{U velocity <ux> <uy> <uz>} - fixed velocity vector.
    \item \codeword{P total <Pa>} - total pressure inlet, the flow follows from the pressure drop.
    \item \codeword{P fixed <Pa>} - static pressure outlet, backflow allowed.
    \item \codeword{T fixed <K>} - inlet temperature, or the temperature of backflow at an outlet. Default \codeword{INIT_T}.
\end{itemize}
Pressures are absolute in Pa. For incompressible runs the setup converts them to kinematic gauge pressure: (p $-$ \codeword{CHARGE_PRESSURE})/$\rho$. The inlet turbulence (k, $\omega$) comes from \codeword{TURB_INTENSITY} (default 0.05) and \codeword{TURB_LENGTH} (\codeword{auto} = 0.07 $\times$ the hydraulic diameter of a round inlet of the same area) at the mean inlet velocity. The setup prints each inlet's area and mean velocity.

Solvers (\codeword{SOLVER} in \codeword{[GLOBAL_SIM_CONTROL]}):
\begin{itemize}
    \item \codeword{rhoSimpleFoam} - steady compressible.
    \item \codeword{buoyantSimpleFoam} - steady compressible with gravity, \codeword{GRAVITY = 0 0 -9.81} (natural draft, hot stacks). Uses \codeword{heRhoThermo} and writes \codeword{constant/g}. The outlet pressure is set with \codeword{prghPressure}, so \codeword{P fixed} is the static pressure at the outlet.
    \item \codeword{simpleFoam} - steady incompressible.
    \item \codeword{rhoPimpleFoam}, \codeword{pimpleFoam} - transient, as for the Stirling cases.
\end{itemize}
Steady runs use \codeword{N_ITER} and \codeword{WRITE_ITER} in \codeword{[GLOBAL_CONTROL]} instead of the cycle keys, under-relax with \codeword{RELAX_P}, \codeword{RELAX_U}, \codeword{RELAX_E}, and stop early once the residuals drop below $10^{-4}$ (p) and $10^{-5}$ (U, e, k, $\omega$). Mesh motion needs a transient solver.

\codeword{postInternal.py} reports, per inlet/outlet group and averaged over the last 50 iterations (steady) or the averaging cycles (transient): \codeword{MDOT_<group>} (kg/s, positive into the domain), \codeword{P_<group>} and \codeword{T_<group>} (area-averaged over the group's patches), \codeword{MDOT_IMBALANCE} (net / inflow) and \codeword{P_DROP} (inlets minus outlets). Steady runs get \codeword{flowConvergence.png}: residuals, flow, p and T against iteration.

\subsection{Fans}
A \codeword{FAN-<name>.stl} planar disk across the duct, slightly larger than the duct so it seals, becomes a fan pressure jump:
\begin{lstlisting}
    [FAN-disk]
    FAN_PRESSURE_RISE = 100
\end{lstlisting}
The rise (Pa) acts in the direction of the disk's surface normal. \codeword{snappyHexMesh} makes an internal face zone. After the zero directory exists, \codeword{runScript} runs \codeword{createBaffles} (\codeword{system/createBafflesDict}), which turns it into the cyclic pair \codeword{<name>_master} / \codeword{<name>_slave} with the \codeword{fan} jump on p (p\_rgh for buoyant runs).

\subsection{Porous Zones}\label{sec:porous}
A closed \codeword{POR-<name>.stl} marks a porous region: the cells inside it become a cell zone with a Darcy--Forchheimer pressure drop, $\nabla p = -(\mu D + \frac{1}{2}\rho |U| F)\,U$ with the superficial velocity U. For a regenerator in a displacer hole, make the solid slightly wider than the hole so its sides are inside the displacer.
\begin{lstlisting}
    [POR-matrix]
    RESISTANCE = wireMesh
    D = 1e7 1e7 1e7
    F = 1e3 1e3 1e3
    E1 = 1 0 0
    E2 = 0 1 0
    POROSITY = 0.9
    WIRE_DIA = 1e-4
    REF_VELOCITY = 1
    COEFFS = local
    DP_DATA =
    DP_LENGTH = 0.1
    MOVES_WITH = none
    HEAT = solved
    HA = auto
    SOLID_T = 300
    SOLID_RHO = 7900
    SOLID_CP = 500
\end{lstlisting}
\begin{itemize}
    \item \codeword{RESISTANCE} - \codeword{coeffs}: \codeword{D} (1/m$^2$) and \codeword{F} (1/m) along the directions \codeword{E1}, \codeword{E2} and their cross product. \codeword{wireMesh}: woven screens (Gedeon--Wood), from \codeword{POROSITY} and \codeword{WIRE_DIA}. \codeword{COEFFS = local} (default): the Reynolds-number dependent parts (the inertial drag and, with \codeword{HA = auto}, the heat transfer) are worked out in every cell and time step from the local gas velocity relative to the matrix; \codeword{REF_VELOCITY} (superficial, m/s) then only sets the values the setup prints. \codeword{COEFFS = ref}: constants taken at \codeword{REF_VELOCITY} for the whole run. \codeword{fit}: least-squares fit of \codeword{DP_DATA} (pairs of superficial velocity in m/s and pressure drop in Pa) measured over \codeword{DP_LENGTH} (m). The setup prints the resulting D and F.
    \item \codeword{HEAT} - \codeword{none}; \codeword{fixed}: the matrix stays at \codeword{SOLID_T} (heater, cooler, heat exchanger core); \codeword{solved}: the matrix temperature changes with its heat capacity (1$-\varepsilon$)$\rho_s c_s$ (\codeword{POROSITY}, \codeword{SOLID_RHO}, \codeword{SOLID_CP}), starting at \codeword{SOLID_T} (regenerator, transient solvers only). \codeword{SOLID_T} takes one value (uniform) or two: the matrix then starts linear from the first value at the $-z$ end of the zone to the second at the $+z$ end (e.g. \codeword{SOLID_T = 600 300} with the hot end at $-z$). A linear start shortens the many cycles a regenerator matrix needs to warm up from a uniform temperature. It is used only when the run has no saved matrix temperature, so restarts keep the current one. The gas receives $h a (T_s - T)$ per unit volume, with \codeword{HA} in W/m$^3$K or \codeword{auto} from the wire-mesh correlation. The term is linearised in the energy, so large values stay stable, and the matrix temperature is updated once per time step with the exact solution of $C_s\,dT_s/dt = h a (T - T_s)$. It is written as the field \codeword{Ts_<zone>} and restarts from it. Compressible runs only.
    \item \codeword{MOVES_WITH} - name of a MOVE body: the porous solid is moved to the start angle with that body before meshing, and its cells then move with the mesh (regenerator in the displacer). The drag then acts on the gas velocity relative to the matrix (below), and the drag force is part of the force on the body.
\end{itemize}
\paragraph{Drag on a moving matrix.} \codeword{explicitPorositySource} works with the absolute velocity, which is right for a fixed matrix only: a matrix moving with the displacer would drag on gas that simply moves along with it. For a zone with \codeword{MOVES_WITH} the setup writes a \codeword{vectorCodedSource} on U instead, with the same Darcy--Forchheimer law on the relative velocity,
\begin{equation}
\mathbf{f} = -\mathbf{C}\cdot(\mathbf{U} - \mathbf{U}_m), \qquad \mathbf{C} = \mu \mathbf{D} + \tfrac{1}{2}\rho\,|\mathbf{U} - \mathbf{U}_m|\,\mathbf{F} ,
\end{equation}
with $\mathbf{D}$, $\mathbf{F}$ the coefficient tensors along \codeword{E1}, \codeword{E2} and their cross product. The matrix velocity $\mathbf{U}_m$ of each cell is the mean motion of its points over the time step (the zone moves rigidly with its body). As in ESI, the isotropic part $\mathrm{tr}(\mathbf{C})/3$ is implicit on the diagonal and the rest is explicit. The force of the gas on the matrix, $\mathbf{C}\cdot(\mathbf{U} - \mathbf{U}_m)$ per unit volume, is stored as \codeword{porDrag_<zone>}; the crank drive adds its integral to the force on the body (\codeword{porousZones} in \codeword{dynamicMeshDict}), and \codeword{postInternal.py} adds it to the body force for prescribed motion.

\paragraph{Drag work.} The drag work $(\mathbf{U} - \mathbf{U}_m)\cdot\mathbf{C}\cdot(\mathbf{U} - \mathbf{U}_m) \ge 0$ (W/m$^3$, $\mathbf{U}_m = 0$ for a fixed zone) heats the gas through the total-energy form of the energy equation (see \codeword{DISSIPATION}); the work the moving matrix does on the gas, $\mathbf{U}_m\cdot\mathbf{C}\cdot(\mathbf{U} - \mathbf{U}_m)$, goes into the gas the same way and is part of the body work. Physically the drag heat appears at the wires; with a large $h a$ the gas passes it on to the matrix within a few time steps, so it makes no practical difference that it is put into the gas. \codeword{DISSIPATION = yes} adds it a second time and logs it as \codeword{porDragHeat_<zone>}.

The heat, drag and drag-work sources are coded fvOptions (\codeword{scalarCodedSource}, \codeword{vectorCodedSource}) in \codeword{constant/fvOptions}, compiled when the solver starts, so a C++ compiler must be available where the job runs. \codeword{postInternal.py} reports \codeword{HEAT_<zone>} (W from the matrix into the gas), \codeword{DRAG_HEAT_<zone>} (W, drag work into the gas), \codeword{TS_MEAN_<zone>}, \codeword{T_MEAN_<zone>} and \codeword{P_MEAN_<zone>}.

\subsubsection{Wire-Screen Regenerator Model}\label{sec:regenModel}
This section gives the equations behind \codeword{RESISTANCE = wireMesh}, \codeword{HA = auto} and \codeword{HEAT = solved}, as they are implemented in \codeword{writeInternal.wireMeshCoeffs} and the heat source in \codeword{constant/fvOptions}.

\paragraph{A stack of screens as one porous zone.}
A regenerator of $N$ stacked woven screens is modelled as one homogeneous porous zone. The correlations below are written for stacks of screens per unit length of matrix, so the number of screens only enters through the length of the zone: the height of the \codeword{POR} solid along the flow is the stack length $L$. For uncompressed plain-weave screens each layer is about two wire diameters thick,
\begin{equation}
L \approx 2 N d_w ,
\end{equation}
with $d_w$ the wire diameter. The porosity of the stack follows from the wire volume per layer. A screen with mesh number $m$ (wires per metre in each direction) holds a wire length of $2m$ per unit area; with cross-section $\pi d_w^2/4$ and layer thickness $2 d_w$, and neglecting the crimp of the weave,
\begin{equation}
1 - \varepsilon = \frac{2m \, \pi d_w^2/4}{2 d_w} = \frac{\pi m d_w}{4},
\qquad \varepsilon \approx 1 - \frac{\pi m d_w}{4} .
\end{equation}
For compressed or sintered stacks the porosity is best taken from the mass $M$ of the stack: $\varepsilon = 1 - M/(\rho_s V_{stack})$.

\paragraph{Geometry of the matrix.}
The wetted surface per unit matrix volume (wires as cylinders, surface/volume $= 4/d_w$ of the solid) and the hydraulic diameter are
\begin{equation}
a = \frac{4(1-\varepsilon)}{d_w}, \qquad d_h = \frac{4\varepsilon}{a} = \frac{\varepsilon d_w}{1-\varepsilon} .
\end{equation}
OpenFOAM works with the superficial velocity $U$ (flow rate over the full cross-section of the zone). The velocity in the pores is $u = U/\varepsilon$, and the Reynolds number of the correlations is built on it:
\begin{equation}
Re = \frac{\rho u d_h}{\mu} = \frac{\rho U d_h}{\varepsilon \mu} .
\end{equation}

\paragraph{Pressure drop.}
Gedeon and Wood (1996, NASA CR-198442, oscillating-flow tests on woven-screen stacks) correlate the friction factor as
\begin{equation}
f = \frac{129}{Re} + 2.91\, Re^{-0.103},
\qquad \frac{\Delta p}{L} = f\, \frac{\rho u^2}{2 d_h} .
\end{equation}
OpenFOAM's \codeword{DarcyForchheimer} model adds the momentum sink
\begin{equation}
\frac{\Delta p}{L} = \mu D\, U + \frac{1}{2}\rho F\, |U|\, U .
\end{equation}
Matching the two term by term, with $u = U/\varepsilon$. The viscous term:
\begin{equation}
\frac{129}{Re}\,\frac{\rho u^2}{2 d_h} = \frac{129\,\mu}{\rho u d_h}\,\frac{\rho u^2}{2 d_h} = \frac{129\,\mu u}{2 d_h^2} = \mu\, \frac{129}{2\varepsilon d_h^2}\, U
\quad\Rightarrow\quad D = \frac{129}{2\varepsilon d_h^2} .
\end{equation}
The inertial term:
\begin{equation}
2.91\,Re^{-0.103}\,\frac{\rho u^2}{2 d_h} = \frac{1}{2}\rho\, \frac{2.91\, Re^{-0.103}}{\varepsilon^2 d_h}\, U^2
\quad\Rightarrow\quad F = \frac{2.91\, Re_{ref}^{-0.103}}{\varepsilon^2 d_h} .
\end{equation}
With \codeword{COEFFS = local} the inertial term uses the Reynolds number of each cell, from its gas density and viscosity and the superficial velocity relative to the matrix $U_r = U - U_{mat}$ ($U_{mat} = 0$ for a static zone): the drag tensor is $\mu D + \frac{1}{2}\rho|U_r|\,F\,(Re/Re_{ref})^{-0.103}$, which is the Gedeon--Wood term at the local $Re$. The setup prints $F$ at \codeword{REF_VELOCITY} ($Re_{ref}$) and a table of $F$, $ha$ and the pressure gradient at 0.1, 1, 5 and 20\,m/s. With \codeword{COEFFS = ref}, $F$ is frozen at $Re_{ref}$; away from it the inertial part is off by $(U/U_{ref})^{-0.103}$, about 7\,\% at twice or half the reference velocity. $D$ is the same along \codeword{E1}, \codeword{E2} and their cross product (isotropic); inside a hole the flow is along the hole anyway.

\paragraph{Heat transfer coefficient.}
Gedeon and Wood give for the same screens
\begin{equation}
Nu = \frac{h\, d_h}{k} = \left(1 + 0.99\,(Re\,Pr)^{0.66}\right) \varepsilon^{1.79} ,
\end{equation}
with the gas conductivity from its viscosity, specific heat and Prandtl number, $k = \mu c_p/Pr$. The volumetric coefficient used by the source is
\begin{equation}
h a = \frac{Nu\, k}{d_h}\cdot\frac{4(1-\varepsilon)}{d_w} ,
\end{equation}
With \codeword{COEFFS = local} it is evaluated per cell and time step, with the cell's $\rho$ and $\mu$ and $Re$ from $|U_r|$ as above. $ha$ grows roughly as $|U_r|^{0.66}$, so in an engine regenerator it changes by about a factor of 7 between flow reversal and peak flow, and a constant value at one velocity gets the matrix exchange wrong over most of the cycle: at 1\,m/s for a regenerator whose flow peaks at about 20\,m/s, the cycle-averaged $ha$ is about half that of the correlation. The setup prints the value at \codeword{REF_VELOCITY}; with \codeword{COEFFS = ref} that value is used throughout. Gas properties for the printed values and for \codeword{COEFFS = ref} are those of \codeword{GAS} at \codeword{CHARGE_PRESSURE} and \codeword{INIT_T} ($\rho$ from the ideal gas law). \codeword{HA} given as a number replaces the correlation with a constant (the drag stays local).

\paragraph{Matrix temperature (local thermal non-equilibrium).}
Each cell of the zone has a gas temperature $T$ and a matrix temperature $T_s$. The matrix exchanges heat only with the gas in the same cell (no conduction along the matrix; stacked screens conduct poorly through their contacts):
\begin{equation}
C_s \frac{dT_s}{dt} = h a\,(T - T_s), \qquad C_s = (1-\varepsilon)\,\rho_s c_s ,
\end{equation}
with $\rho_s$ = \codeword{SOLID_RHO}, $c_s$ = \codeword{SOLID_CP}. Gas and matrix are stepped together, implicitly (backward Euler), so that the heat the gas receives in a time step is exactly the heat the matrix gives up:
\begin{equation}
C_s\,\frac{T_s^{n+1} - T_s^{n}}{\Delta t} = h a\,(T^{n+1} - T_s^{n+1}) .
\end{equation}
Eliminating $T_s^{n+1}$, the heat into the gas is
\begin{equation}
Q_s = h a\,(T_s^{n+1} - T^{n+1}) = h_{eff}\,(T_s^{n} - T^{n+1}), \qquad h_{eff} = \frac{h a\, C_s}{C_s + h a\,\Delta t} ,
\end{equation}
which is stable for any time step. After the energy equation is solved, $Q_s$ is evaluated with the new gas temperature and the matrix is updated with $T_s^{n+1} = T_s^{n} - Q_s\,\Delta t/C_s$, so the two sides balance to round-off. With \codeword{HEAT = fixed}, $C_s = 0$: $h_{eff} = h a$ and $T_s$ stays at \codeword{SOLID_T}.

\paragraph{Heat source in the gas energy equation.}
The gas receives $Q_s = h_{eff}\,(T_s^{n} - T)$ per unit volume. For a large $h a$ an explicit source makes the energy equation stiff, so it is linearised in the internal energy $e$: with $T \approx T^* + (e - e^*)/c_v$ around the current state $(T^*, e^*)$,
\begin{equation}
Q_s \approx \underbrace{h_{eff}\,(T_s^{n} - T^*) + \frac{h_{eff}}{c_v}\, e^*}_{\textrm{explicit}} \;-\; \underbrace{\frac{h_{eff}}{c_v}\, e}_{\textrm{implicit}} ,
\end{equation}
with $c_v = c_p - R_u/M$ of the gas. The implicit part goes on the matrix diagonal, the explicit part into the source; with $c_v$ constant the linearisation is exact. The heat from the matrix into the gas, $\int_V Q_s\,dV$, is written by the \codeword{porHeatToGas_<zone>} function object (negative while the gas heats the matrix) and reported as \codeword{HEAT_<zone>}.

\paragraph{Fit to measured data (\texttt{RESISTANCE = fit}).}
For measured pairs $(U_i, \Delta p_i)$ over a sample of length \codeword{DP_LENGTH} $= L$, $D$ and $F$ are the least-squares solution through the origin of
\begin{equation}
\left(\begin{array}{cc} \mu U_1 & \frac{1}{2}\rho U_1|U_1| \\ \vdots & \vdots \\ \mu U_n & \frac{1}{2}\rho U_n|U_n| \end{array}\right)
\left(\begin{array}{c} D \\ F \end{array}\right) =
\left(\begin{array}{c} \Delta p_1/L \\ \vdots \\ \Delta p_n/L \end{array}\right) .
\end{equation}
The setup prints the largest misfit and stops if either coefficient comes out negative.

\paragraph{Known modelling errors.}
\begin{itemize}
    \item The solver treats the whole zone as gas volume (no $\varepsilon$ in the time derivative), so the gas volume and gas heat capacity in the matrix are over-counted by $(1-\varepsilon) V_{zone}$. Count it as extra dead volume.
    \item With \codeword{COEFFS = ref}, $F$ and $h a$ use a fixed $Re_{ref}$ (see above); the default \codeword{COEFFS = local} evaluates them per cell.
    \item The zone is meshed conformally: snappy selects the cells inside the \codeword{POR} solid and then snaps the zone boundary onto the surface, so the stack length is the drawn one to snapping accuracy (on the test pipe the zone volume matched the drawn cylinder within 0.5\,\%). Flat end faces and five to ten cells along the stack snap cleanly and resolve the temperature along it.
    \item No axial conduction in the matrix.
\end{itemize}

\paragraph{Worked example.}
Fifty 100-mesh screens, $d_w = 0.1$\,mm, stainless steel ($\rho_s = 7900$, $c_s = 500$), air at 1\,bar and 300\,K ($\rho = 1.161$\,kg/m$^3$, $\mu = 1.81\cdot10^{-5}$\,Pa\,s, $c_p = 1005$\,J/kg\,K, $Pr = 0.71$), \codeword{REF_VELOCITY} $= 1$\,m/s:
\begin{itemize}
    \item $m = 100/0.0254 = 3937$\,m$^{-1}$, $\varepsilon = 1 - \pi \cdot 3937 \cdot 10^{-4}/4 = 0.691$, $L \approx 2 \cdot 50 \cdot 0.1$\,mm $= 10$\,mm.
    \item $d_h = 0.691 \cdot 10^{-4}/0.309 = 2.234\cdot10^{-4}$\,m, $a = 1.237\cdot10^{4}$\,m$^{-1}$, $u = 1.448$\,m/s, $Re = 20.7$.
    \item $D = 129/(2 \cdot 0.691 \cdot (2.234\cdot10^{-4})^2) = 1.87\cdot10^{9}$\,m$^{-2}$, $Re^{-0.103} = 0.732$, $F = 2.91 \cdot 0.732/(0.691^2 \cdot 2.234\cdot10^{-4}) = 2.00\cdot10^{4}$\,m$^{-1}$.
    \item At 1\,m/s: $\mu D U = 33.9$\,kPa/m, $\frac{1}{2}\rho F U^2 = 11.6$\,kPa/m, together 45.5\,kPa/m, or 455\,Pa over the 10\,mm stack.
    \item $Re\,Pr = 14.7$, $Nu = (1 + 0.99 \cdot 14.7^{0.66}) \cdot 0.691^{1.79} = 3.53$, $k = 0.0256$\,W/m\,K, $h = 405$\,W/m$^2$K, $h a = 5.0\cdot10^{6}$\,W/m$^3$K.
    \item $C_s = 0.309 \cdot 7900 \cdot 500 = 1.22\cdot10^{6}$\,J/m$^3$K, so the matrix follows the gas with the time constant $C_s/(h a) = 0.24$\,s.
\end{itemize}
These are the values the setup prints for \codeword{POROSITY = 0.6908}, \codeword{WIRE_DIA = 1e-4}, \codeword{REF_VELOCITY = 1}.

\subsubsection{Random-Fibre Model (Steel Wool)}\label{sec:fiberModel}
\codeword{RESISTANCE = randomFiber} models a loose tangle of thin fibres, such as steel wool stuffed into the hot end space. The idea is simple: a flat hot plate only touches the gas next to it, but wool spreads a huge metal surface through the whole space, so the gas picks up heat much faster. The wool also gets squashed and released as the displacer moves.

The model does four things in every cell, every time step:
\begin{enumerate}
    \item keeps track of how much fibre is in the cell,
    \item slows the gas down (drag),
    \item swaps heat between the gas and the fibres,
    \item spreads heat through the wool itself, and from the hot plate into the wool.
\end{enumerate}

\begin{lstlisting}
    [POR-wool]
    RESISTANCE = randomFiber
    WIRE_DIA = 5e-5
    SOLID_MASS = 0.05
    SOLID_RHO = 7850
    SOLID_CP = 470
    SOLID_K = 50
    RADIATION = yes
    WALL_CONTACT_PIDS = hot
    WALL_CONTACT = 1
    MOVES_WITH = mesh
    HEAT = solved
    HA = auto
    SOLID_T = 600
    REF_VELOCITY = 0.5
\end{lstlisting}
\begin{itemize}
    \item \codeword{WIRE_DIA} - fibre thickness (m).
    \item \codeword{SOLID_MASS} - how much wool is in the space (kg). Or give \codeword{POROSITY} for a fixed, evenly packed bed. Wool that moves with the mesh needs \codeword{SOLID_MASS}.
    \item \codeword{SOLID_RHO}, \codeword{SOLID_CP}, \codeword{SOLID_K} - density, specific heat and conductivity of the fibre metal. \codeword{SOLID_K} defaults to 50\,W/m\,K (ordinary steel); stainless steel is about 15.
    \item \codeword{RADIATION} - \codeword{yes} lets heat also travel through the wool as radiation.
    \item \codeword{WALL_CONTACT_PIDS} - the walls the wool touches, e.g. the hot plate. \codeword{WALL_CONTACT} (0 to 1) says how good that contact is; 1 = perfect.
    \item \codeword{MOVES_WITH = mesh} - the wool is squashed and stretched with the mesh. Draw the \codeword{POR} solid as a box filling the end space.
    \item \codeword{HEAT = solved}, \codeword{HA = auto}, one \codeword{SOLID_T} (start temperature of the wool).
    \item \codeword{REF_VELOCITY} - only used for the table the setup prints; the run itself uses the real local gas speed.
\end{itemize}

\paragraph{1. How much fibre is in each cell.}
Each cell stores its fibre fraction $\varphi$: the share of its volume that is metal (field \codeword{solidFrac_<zone>}); the rest, $\varepsilon = 1 - \varphi$, is gas. The fibres stay in their cell, so when the displacer squashes a cell to half its size, $\varphi$ doubles. At every (re)start the total is reset to exactly \codeword{SOLID_MASS}, so no wool is gained or lost, also not when the mesh is remade. The fibres move with the mesh, and everything below uses the gas speed \emph{relative to the fibres}.

\paragraph{2. Drag.}
Gas flowing through the wool is slowed down, more so the denser the wool and the faster the flow. No measured data exist for wool as loose as steel wool (95--99\,\% empty), so two textbook cases are combined:
\begin{itemize}
    \item \textbf{Dense metal felt} - correlations measured on packed fibre mats by Gedeon and Wood (NASA CR-198442), valid up to about 84\,\% empty.
    \item \textbf{Single fibres} - each fibre treated as a thin wire in the flow, on its own (wire drag from White, and from Lamb at very low speed). This is the right picture when the fibres are far apart.
\end{itemize}
The model uses the felt data for dense wool, the single-wire picture for loose wool (above 97\,\% empty), and blends smoothly between the two. So as the displacer squashes the wool, the model moves from one to the other. The drag force on the fibres is logged as \codeword{porDrag_<zone>}.

\paragraph{3. Heat between gas and fibres.}
How fast heat moves between gas and fibres depends on the total fibre surface ($4\varphi/d_w$ per m$^3$: thinner and denser wool has more) and on how fast the gas flows past it. The same two cases are used and blended in the same way (felt: Gedeon and Wood; single wires: Churchill and Bernstein). The heat that leaves the gas is exactly the heat that enters the fibres, and vice versa, so no energy is lost. The fibre temperature is the field \codeword{Ts_<zone>}; the heat from the fibres into the gas is logged as \codeword{porHeatToGas_<zone>}.

\paragraph{4. Heat through the wool and from the hot plate.}
This step lets the wool work like the fins of a radiator. Without it, the wool could only take heat from the gas and never directly from the hot plate. Heat travels through the wool in two ways, added together:
\begin{equation}
k_{eff} = \underbrace{\frac{k_s\,\varphi}{3}}_{\textrm{along the fibres}} + \underbrace{\frac{16\,\sigma\,T_s^3\, d_w}{3\,\varphi}}_{\textrm{radiation}} .
\end{equation}
\begin{itemize}
    \item \textbf{Along the fibres.} Only the metal conducts, and only along each fibre. A fibre pointing in the direction of the heat flow carries it fully, a fibre lying across carries none. Averaged over fibres pointing every which way, one third of the metal is effectively in line, hence $\varphi k_s/3$. This assumes the fibres are long and touch each other perfectly, so it is the \emph{best case}. Real wool is crinkled and the fibres only touch at points, so the true value is lower. Lower \codeword{SOLID_K} if you know the wool conducts worse.
    \item \textbf{Radiation.} Hot fibres glow (in infrared) and pass heat on to their neighbours. The formula is the standard one for a dense, cloudy material (Rosseland). It works when a ray of light can only get a few millimetres before hitting a fibre. In wool squashed very flat this is no longer quite true, and the formula then overestimates the radiation.
\end{itemize}
For scale: at 600\,K, wool of 50\,$\mu$m steel conducts about 0.5 to 0.9\,W/m\,K, ten to twenty times better than still air (0.045\,W/m\,K).

Wool cells that touch a \codeword{WALL_CONTACT_PIDS} wall also take heat straight from it, scaled by \codeword{WALL_CONTACT}. Nobody can predict how well wool touches a plate, so treat \codeword{WALL_CONTACT} as a knob: run with a lower and a higher value to see how much it matters. This heat is logged as \codeword{porWallToFiber_<zone>}. Count it as heat input in energy balances.

\paragraph{What to keep in mind.}
\begin{itemize}
    \item The drag and heat-transfer numbers are uncertain for loose steel wool: there is no measured data at 95--99\,\% empty, and the two textbook cases differ by 4 to 10 times there. Measurements on the actual wool would fix this.
    \item Heat conduction through the wool is a best case; radiation is rough when the wool is squashed flat.
    \item How well the wool touches the hot plate (\codeword{WALL_CONTACT}) is a guess you supply.
    \item Not modelled: the springiness of the wool pushing on the displacer, gaps between wool and walls, uneven packing (the wool starts evenly spread), and the small volume the fibres take away from the gas.
\end{itemize}

\subsection{Function Objects}
Written to \codeword{system/internalFunctionObjects}, every time step:
\begin{itemize}
    \item \codeword{gasAvg} - volume-averaged p (and T, $\rho$ for compressible).
    \item \codeword{gasMass} - total gas mass (compressible). It must stay constant in a closed volume. A drift shows a mesh motion or solver tolerance problem.
    \item \codeword{wallHeatFlux} - heat flux on every wall patch (compressible), positive = into the gas.
    \item \codeword{force_<body>} - force on the moving faces of each body.
    \item \codeword{p_<patch>} - area-averaged pressure on each moving zone face.
    \item \codeword{por_<zone>} - volume averages over each porous zone (p, T, matrix temperature).
    \item \codeword{porHeatToGas_<zone>} - heat from the matrix into the gas (W, \codeword{HEAT} fixed or solved).
    \item \codeword{porDragHeat_<zone>} - drag work added to the gas (W, only with \codeword{DISSIPATION = yes}).
    \item \codeword{porDrag_<zone>} - force of the gas on a porous zone moving with a body (N, vector).
    \item \codeword{porWallToFiber_<zone>} - heat from the contact walls into a random-fibre bed (W, section~\ref{sec:fiberModel}).
    \item \codeword{fiber_<zone>} - fibre volume of a random-fibre zone (m$^3$; constant when \codeword{SOLID_MASS} is kept).
\end{itemize}

\subsection{Post-Processing}\label{sec:internalPost}
\codeword{postScript} (generated by the setup, chained after \codeword{runScript} by \codeword{submitScript}, or run locally with \codeword{./postScript}) runs, in order:
\begin{enumerate}
    \item \codeword{postInternal.py} - the summary \codeword{summary.csv} and the plots \codeword{pV.png} and \codeword{cycleConvergence.png}, all in the case folder.
    \item \codeword{gsheetInternal.py} - the summary row in the database sheet (\codeword{GSHEET = True}).
    \item \codeword{pvPostInternal.py} (pvbatch) - one image per written time and panel, \codeword{postProcessing/images/<panel>/}.
    \item \codeword{imagesToMovies.py} (ffmpeg) - one video per panel, \codeword{postProcessing/movies/<case>_<panel>.mp4}.
\end{enumerate}
Logs: \codeword{log.postInternal}, \codeword{log.gsheetInternal}, \codeword{log.pvPostInternal}, \codeword{log.imagesToMovies}.

\subsubsection{Summary}
\codeword{postUtilities/postInternal.py} writes one summary per case to \codeword{summary.csv} in the case folder (key, value), and the plots \codeword{pV.png} (p--V loops of the last \codeword{--pvCycles} cycles, default 3, compressible runs) and \codeword{cycleConvergence.png} (per cycle: work per body, mean gas T, gas mass drift, heat per BC group, crank speed); \codeword{--noPlots} skips them:
\begin{lstlisting}
    python3 /path/to/caseSetup/postUtilities/postInternal.py [cases ...] [--avgCycles N] [--pvCycles N] [--noPlots] [--gsheet]
\end{lstlisting}
\begin{itemize}
    \item Run information: job, name, description, run and completion dates, cell count, mesh, solve and post times (hours), the main setup values and the BC groups.
    \item Status: end time, cycles done, finished.
    \item Results, averaged over the last \codeword{N} complete cycles: work per body per cycle, \codeword{AVG_ENERGY} (J/cycle, positive = work out of the gas; a displacer alone gives a negative value, the work needed to drive it), \codeword{AVG_POWER}, mean/max/min gas pressure, mean T and $\rho$, pressure mean and amplitude on each moving face, force amplitude per body.
    \item Compressible: gas mass drift, p--V work, gas volume range, heat per BC group (W, positive into the gas).
    \item Crank drive: work from $\oint T\,d\theta$, cycles counted by crank angle, mean/min/max speed, speed change over the last cycle (convergence), mean gas torque, load power and electrical power.
\end{itemize}
Restarted runs are handled: the function object directories are merged in time order. The equations behind every value are in Section~\ref{sec:postMath}.

\codeword{postUtilities/gsheetInternal.py} pushes the summaries to the database sheet, one row per case found by job and name (\codeword{--comment}, \codeword{--dryRun}). The sheet id and the service account key are read from \codeword{~/.config/caseSetup/gsheet.ini} or the \codeword{INTERNAL_GSHEET_*} environment variables. Keep the key out of git.

\subsubsection{Images}
\codeword{postUtilities/pvPostInternal.py} is set up like \codeword{pvPost}: the file \codeword{pvPostInternalSetup} in the case directory selects the panels, and a json variable dict defines the variables. The setup copies the default file into the case the first time and keeps an existing copy. Keys missing from the case file take the default values.
\begin{lstlisting}
    [PV_POST_MAIN]
    VARIABLE_DICT = default
    RESOLUTION = 1920 1080
    EVERY = 1
    LAST_CYCLES = all
    SLICE_IMG = True
    SURFACE_IMG = True
    VOLUME_IMG = True
    CLIP_IMG = True
    CREATE_SCALAR_BAR = True

    [SLICE]
    NAMES = yMid xMid zMid
    NORMALS = Y X Z
    ORIGINS = center center center
    FOLLOW = none none none
    VARIABLES = Umag Uz p T

    [SURFACE]
    PATCHES = MOVE
    VIEWS = Iso
    VARIABLES = p T

    [VOLUME]
    VIEWS = Iso
    VARIABLES = T Umag
    OPACITY = 0.8
    OPACITY_RAMP = auto
    UNIT_DISTANCE = auto
    CONTEXT_PATCHES = MOVE
    CONTEXT_OPACITY = 0.15
    RESAMPLE = 250
    OUTLINE = False

    [CLIP]
    NAMES = yClip
    NORMALS = -Y
    ORIGINS = center
    VARIABLES = T Umag p

    [MOVIES]
    ENABLE = True
    FPS = 10
\end{lstlisting}
\begin{itemize}
    \item \codeword{EVERY} - every N-th written time (counted back from the last one, so the last time is always used). \codeword{LAST_CYCLES} - \codeword{all}, or only the last N cycles. Both can be given on the command line as \codeword{--every} and \codeword{--lastCycles}.
    \item \codeword{RESOLUTION} - image size in pixels, also the video size.
    \item \codeword{[SLICE]} - one entry per slice in each of the lists \codeword{NAMES}, \codeword{NORMALS} (\codeword{X}, \codeword{Y}, \codeword{Z}), \codeword{ORIGINS} and \codeword{FOLLOW}. Origin \codeword{center} = middle of the mesh box, \codeword{loc} = \codeword{LOC_IN_MESH}, or \codeword{x,y,z}. \codeword{FOLLOW = <body>} moves the slice in z with that body (from the motion table, or with the crank drive from the crank log and the crank table), so a z slice through the displacer holes stays in the holes. Every slice is drawn for every variable in \codeword{VARIABLES}.
    \item \codeword{[SURFACE]} - the patches whose name contains one of the \codeword{PATCHES} keywords (e.g. \codeword{MOVE}: the displacer), seen from each of the \codeword{VIEWS} (\codeword{Iso}, \codeword{Top}, \codeword{Bottom}, \codeword{Front}, \codeword{Back}, \codeword{Left}, \codeword{Right}).
    \item \codeword{[VOLUME]} - volume rendering of the gas, seen from each of the \codeword{VIEWS}: semi-transparent colour through the whole gas volume, so the hot and cold gas and the jets from the holes show in 3D. The gas is resampled to a regular grid with \codeword{RESAMPLE} points along the longest side of the mesh box (ParaView's volume rendering of the unstructured mesh skips the polyhedral cells of snappyHexMesh). The opacity rises over the colour range up to \codeword{OPACITY}: \codeword{OPACITY_RAMP} \codeword{rising} (clear at the low end, so hot gas or fast flow stands out), \codeword{falling}, \codeword{ends} (clear in the middle, for signed variables like \codeword{Uz}), or \codeword{auto} (\codeword{ends} for symmetric ranges, else \codeword{rising}). \codeword{UNIT_DISTANCE} is the depth over which that opacity builds up: \codeword{auto} = 0.05 $\times$ the shortest side of the mesh box, larger values are more transparent. The patches matching \codeword{CONTEXT_PATCHES} (blank = none) are drawn in grey with \codeword{CONTEXT_OPACITY}. \codeword{OUTLINE = True} also draws the edges of the mesh box.
    \item \codeword{[CLIP]} - half of the gas volume, cut by a plane through \codeword{ORIGINS} (as for slices) with normal \codeword{NORMALS} (\codeword{X}, \codeword{Y}, \codeword{Z}, \codeword{-X}, \codeword{-Y}, \codeword{-Z}). The normal points to the removed half, and the camera looks at the cut from that side, tilted so the depth shows: the cut face and the gas boundary (walls, displacer faces, bore) are coloured by each variable in \codeword{VARIABLES}.
    \item \codeword{VOLUME_IMG}, \codeword{CLIP_IMG} (and \codeword{SLICE_IMG}, \codeword{SURFACE_IMG}) switch the panel types on or off. \codeword{--panels volume,clip} on the command line makes only the listed types.
    \item \codeword{VARIABLE_DICT} - \codeword{default} (\codeword{postUtilities/default/defaultInternalVariables}) or the path to your own copy. Each variable has an \codeword{equation} (ParaView calculator, \codeword{PREF} is replaced by the charge pressure), a \codeword{label}, a \codeword{range} and a ParaView colour map. Slice variables (also used by the volume and clip panels): \codeword{Umag}, \codeword{Ux}, \codeword{Uy}, \codeword{Uz}, \codeword{p}, \codeword{pGauge} (p minus the charge pressure), \codeword{T}, \codeword{rho}. Surface variables: \codeword{p}, \codeword{pGauge}, \codeword{T}, \codeword{wallHeatFlux}. Variables whose fields are not in the case (e.g. \codeword{T} for \codeword{pimpleFoam}) are skipped.
    \item \codeword{range} - \codeword{auto}: minimum to maximum over all selected times, the same on every slice. \codeword{symmetric}: $\pm$ the largest absolute value. \codeword{[min, max]}: fixed. The range is fixed over the run, so the colours mean the same in every frame of a video.
\end{itemize}
Each image has a title with the case, panel, variable, time and crank angle. Run it by hand from the case directory:
\begin{lstlisting}
    pvbatch --force-offscreen-rendering /path/to/caseSetup/postUtilities/pvPostInternal.py [--every N] [--lastCycles N] [--panels slice,surface,volume,clip]
\end{lstlisting}

\subsubsection{Videos}
\codeword{postUtilities/imagesToMovies.py} makes one mp4 (H.264) per image folder, one image per frame at \codeword{[MOVIES] FPS} (default 10; \codeword{--fps} overrides it). \codeword{ENABLE = False} turns the videos off (\codeword{--force} makes them anyway). It needs \codeword{ffmpeg} on the path.
\begin{lstlisting}
    python3 /path/to/caseSetup/postUtilities/imagesToMovies.py [--fps N]
\end{lstlisting}

\subsubsection{Copy Back to Google Drive}
\codeword{postUtilities/copyBackInternal.py} uploads the run data of a case to \codeword{<folder>/<JOB>/<case>/} on Google Drive, keeping the case's own sub-paths: logs and Slurm \codeword{.out} files, \codeword{internalSetup}, \codeword{fullCaseSetupDict}, \codeword{pvPostInternalSetup}, \codeword{summary.csv} and the plots, every \codeword{postProcessing/**/*.dat}, the images and the videos. Only new or changed files are uploaded (MD5 compared with the copy on Drive). It is the last step of \codeword{postScript} when \codeword{COPY_BACK = True} (\codeword{[GLOBAL_COMPUTE_SETUP]}), and can be run by hand:
\begin{lstlisting}
    python3 /path/to/caseSetup/postUtilities/copyBackInternal.py [cases ...] [--dryRun] [--noImages]
\end{lstlisting}
Settings in \codeword{~/.config/caseSetup/drive.ini} (or \codeword{--config}): a \codeword{[drive]} section with \codeword{folderID} (the id or the folder URL) and \codeword{credentials} (the service account key, default \codeword{~/.config/caseSetup/internal_cfd_credentials.json}); \codeword{INTERNAL_DRIVE_FOLDER} and \codeword{INTERNAL_DRIVE_CREDENTIALS} override it. The folder must be on a Shared Drive, because service accounts have no storage of their own, and shared with the service account as Content manager. Needs \codeword{google-api-python-client} and \codeword{google-auth}.

\subsection{Post-Processing and Motion Calculations}\label{sec:postMath}
This section gives the equations behind the numbers in \codeword{summary.csv}, \codeword{pV.png}, \codeword{cycleConvergence.png} and \codeword{flowConvergence.png} (\codeword{postUtilities/postInternal.py}), the crank drive (\codeword{foamLibs/stirlingMotion}) and the mesh motion report printed by the setup (\codeword{motionUtilities.py}). Units are SI throughout, angles in the equations are in radians unless given in degrees.

\subsubsection{Input Data and Sign Conventions}
\begin{itemize}
    \item The function objects write one row per time step to \codeword{postProcessing/<name>/<startTime>/*.dat}. After a restart there are several start times. They are merged in time order, and a later file replaces the overlapping times of an earlier one.
    \item $\mathbf{S}_f$ is the face area vector of a boundary face. OpenFOAM points it out of the gas, so on a moving body it points into the body. The force of the gas on a body, $\mathbf{F} = \sum_f p\,\mathbf{S}_f$ (plus the viscous part), is then positive in the direction the gas pushes the body.
    \item Work and power are positive when they leave the gas (the gas pushes a body the way it moves: engine output). Heat is positive into the gas. Mass flow is positive into the domain.
    \item The motion tables give the displacement $\Delta z(t)$ of each body from the meshed position. The solver interpolates them linearly in time.
\end{itemize}

\subsubsection{Cycles and Time Averages}
\paragraph{Cycle windows.}
With a prescribed motion the run starts at $t = 0$ at \codeword{START_ANGLE}, and cycle $k$ ($k = 1, 2, \ldots$) runs from $t_{k-1}$ to $t_k$ with
\begin{equation}
t_k = \frac{k}{f}, \qquad N_{done} = \left\lfloor t_{end} f \right\rfloor ,
\end{equation}
$f$ = \codeword{FREQ}, $t_{end}$ = the last time in \codeword{gasAvg}. Only complete cycles count. The results are averaged over the last $N = \min($\codeword{--avgCycles}$, N_{done})$ cycles, the window
\begin{equation}
t_0 = \frac{N_{done} - N}{f}, \qquad t_1 = \frac{N_{done}}{f} .
\end{equation}
With the crank drive the speed is not fixed, so cycles are counted by the crank angle $\theta$ (accumulated, degrees, from the log \codeword{crank.dat}) instead of time: $N_{done} = \lfloor(\theta_{end} - \theta_0)/360\rfloor$ with $\theta_0$ the start angle, and the cycle edges are the times at which $\theta = \theta_0 + 360k$, interpolated linearly in the log. If the crank stops or turns backwards, only the cycles before the first stop are used.

\paragraph{Time average.}
The time step varies, so a plain average of the samples would weigh the parts of the cycle with small time steps too much. Every average is a time average with the trapezoidal rule, with the window ends interpolated:
\begin{equation}
\overline{y} = \frac{1}{t_1 - t_0}\int_{t_0}^{t_1} y\,dt \approx \frac{1}{t_1 - t_0}\sum_i \frac{y_i + y_{i+1}}{2}\,(t_{i+1} - t_i) .
\end{equation}
Example: samples $y = 2, 4, 4$ at $t = 0, 1, 3$ give $\overline{y} = (3 \cdot 1 + 4 \cdot 2)/3 = 3.67$, the sample mean would be 3.33. All \codeword{_MEAN} values are computed like this. \codeword{_MAX} and \codeword{_MIN} are the extreme samples in the window, and an amplitude is half the range:
\begin{equation}
y_{amp} = \frac{1}{2}\left(\max_{t_0 \le t \le t_1} y - \min_{t_0 \le t \le t_1} y\right) .
\end{equation}
This is used for \codeword{P_<patch>_AMP} (area-averaged pressure on each moving face) and \codeword{FORCE_Z_AMP_<body>}.

\subsubsection{Work and Power, Prescribed Motion}
The work done by the gas on a body over the window, per cycle:
\begin{equation}
W_{body} = \frac{1}{N}\int_{t_0}^{t_1} F_z(t)\, v(t)\, dt ,
\end{equation}
with $F_z$ the \codeword{total_z} column of \codeword{force_<body>} (pressure and viscous), plus the $z$ component of \codeword{porDrag_<zone>} for porous zones moving with the body (reported separately as \codeword{WORK_DRAG_<body>}), and $v$ the body velocity from its motion table. Because the solver interpolates the table linearly in time, the mesh moves with a constant velocity between two table points $t_j$ and $t_{j+1}$:
\begin{equation}
v(t) = \frac{\Delta z_{j+1} - \Delta z_j}{t_{j+1} - t_j}, \qquad t_j \le t < t_{j+1} .
\end{equation}
The same piecewise constant velocity is used in the integral, so the work matches what the mesh actually did. Then
\begin{equation}
\mathrm{AVG\_ENERGY} = \sum_{bodies} W_{body}, \qquad \mathrm{AVG\_POWER} = \mathrm{AVG\_ENERGY}\cdot f .
\end{equation}
A body that only moves gas around (a displacer alone) has $W < 0$: it takes work to drive it, the flow loss through the holes and the gap.

The force function objects use the absolute pressure. A constant pressure $p_c$ on a face of area $A$ adds $p_c A \oint dz = 0$ over a closed cycle, so the work does not depend on the pressure level; the force amplitude does not either.

\subsubsection{Crank Drive}
\paragraph{In the solver, every time step.}
The force of the gas on each body along the axis, pressure only, relative to the pressure behind the body $p_{ref}$ (\codeword{P_BUFFER}):
\begin{equation}
F_{body} = \sum_{f \in body} (p_f - p_{ref})\,\mathbf{S}_f \cdot \mathbf{e}_z .
\end{equation}
Porous zones moving with the body add the gas drag on them, $\sum_{cells} (\mathbf{f}_{drag}\cdot\mathbf{e}_z)\,V$ from \codeword{porDrag_<zone>} of the previous step. For \codeword{pimpleFoam} $p$ is kinematic, so $F$ is multiplied by the gas density at the charge state and $p_{ref} = 0$. The viscous force is left out; in a typical displacer run it is about $10^{-6}$ of the pressure force.

The torque on the crank follows from the virtual work, $F\,dz = T\,d\theta$:
\begin{equation}
T_{body} = F_{body}\,\frac{dz_{body}}{d\theta}, \qquad \frac{dz}{d\theta} \approx \frac{z(\theta + h) - z(\theta - h)}{2h}, \quad h = 0.25^\circ ,
\end{equation}
with $z(\theta)$ from the crank table \codeword{<body>Crank} (in m per rad). The flywheel:
\begin{equation}
I\,\frac{d\omega}{dt} = T_{gas} - T_0 - c\,\omega, \qquad T_{gas} = \sum_{bodies} T_{body} ,
\end{equation}
with $I$ = \codeword{INERTIA}. The load is split into a constant part $T_0$ and a speed-proportional part $c\,\omega$:
\begin{center}
\begin{tabular}{l l l}
\hline
\codeword{LOAD} & $T_0$ & $c$ \\
\hline
\codeword{torque} & $T_{fr}\,\mathrm{sign}(\omega) + T_{load}$ & $c_{fr}$ \\
\codeword{damper} & $T_{fr}\,\mathrm{sign}(\omega)$ & $c_{fr} + c_{damp}$ \\
\codeword{generator} & $T_{fr}\,\mathrm{sign}(\omega)$ & $c_{fr} + k^2/(R_{load} + R_{int})$ \\
\hline
\end{tabular}
\end{center}
$T_{fr}$ = \codeword{FRICTION_TORQUE}, $c_{fr}$ = \codeword{FRICTION_COEFF}, $T_{load}$ = \codeword{LOAD_TORQUE}, $c_{damp}$ = \codeword{LOAD_DAMPING}, $k$ = \codeword{GEN_K}. The generator line comes from the electrical circuit: the back-EMF $k\omega$ drives the current $I_e = k\omega/(R_{load} + R_{int})$, the generator torque is $k I_e = k^2\omega/(R_{load} + R_{int})$, and the power in the load resistor is
\begin{equation}
P_{elec} = I_e^2 R_{load} .
\end{equation}
The gas torque is explicit (from the pressure of the previous step), the speed-proportional part implicit, so a large damping cannot make the step unstable:
\begin{equation}
\omega^{n+1} = \frac{I\,\omega^n + \Delta t\,(T_{gas} - T_0)}{I + \Delta t\, c}, \qquad
\theta^{n+1} = \theta^n + \frac{\Delta t}{2}\left(\omega^n + \omega^{n+1}\right) .
\end{equation}
The logged load torque is $T_0 + c\,\omega^{n+1}$, so it includes the friction. With \codeword{LOAD = speed}, $\omega$ stays at $2\pi f$ and the load torque is set equal to the gas torque. With mesh motion outer correctors the step is redone from $(\theta^n, \omega^n)$ with the newer pressure.

\paragraph{In the summary.}
The work per body is the torque integrated over the crank angle,
\begin{equation}
W_{body} = \frac{1}{N}\int_{\theta(t_0)}^{\theta(t_1)} T_{body}\, d\theta ,
\end{equation}
which is the same as $\int F v\,dt$ because $v = (dz/d\theta)\,\omega$. Speed and power:
\begin{equation}
\mathrm{FREQ\_MEAN} = \frac{N}{t_1 - t_0}, \qquad \mathrm{AVG\_POWER} = \mathrm{AVG\_ENERGY}\cdot \mathrm{FREQ\_MEAN},
\end{equation}
\begin{equation}
\mathrm{FREQ\_MIN/MAX} = \frac{\omega_{min/max}}{2\pi} ,
\end{equation}
\begin{equation}
\mathrm{TORQUE\_GAS\_MEAN} = \frac{1}{\theta(t_1) - \theta(t_0)}\int T_{gas}\,d\theta ,
\end{equation}
\begin{equation}
\mathrm{POWER\_LOAD} = \overline{T_{load}\,\omega}, \qquad \mathrm{POWER\_ELEC} = \overline{P_{elec}} .
\end{equation}
\codeword{POWER_LOAD} includes the friction losses. The speed change of the last cycle is a convergence measure:
\begin{equation}
\mathrm{FREQ\_CHANGE} = \frac{1}{t_{N_{done}} - t_{N_{done}-1}} - \frac{1}{t_{N_{done}-1} - t_{N_{done}-2}} .
\end{equation}
At a periodic state the mean gas torque equals the mean load torque, $\overline{T_{gas}} = \overline{T_0 + c\,\omega}$, and \codeword{FREQ_CHANGE} goes to zero.

\subsubsection{Gas State, p--V Loop and Mass}
\paragraph{Gas volume.}
\codeword{gasAvg} writes the volume averages $\bar{p}$, $\bar{T}$, $\bar{\rho}$ and \codeword{gasMass} the gas mass $m = \int \rho\, dV$. Since $\bar{\rho} = m/V$ by definition, the gas volume follows exactly, without the mesh volume:
\begin{equation}
V(t) = \frac{m(t)}{\bar{\rho}(t)} .
\end{equation}
\codeword{GAS_VOLUME_MIN} and \codeword{_MAX} are its extremes in the window. \codeword{P_MEAN}, \codeword{T_MEAN}, \codeword{RHO_MEAN} are the time averages of the volume averages.

\paragraph{p--V work.}
\begin{equation}
\mathrm{PV\_WORK} = \frac{1}{N}\oint \bar{p}\, dV \approx \frac{1}{N}\sum_i \frac{\bar{p}_i + \bar{p}_{i+1}}{2}\,(V_{i+1} - V_i)
\end{equation}
over the window. A loop run clockwise in the p--V plot gives positive work (out of the gas). \codeword{PV_WORK} uses one mean pressure for the whole gas. The body work uses the pressure on each face, so the two differ by the pressure differences inside the engine (flow loss through the regenerator and the holes): $\mathrm{PV\_WORK} - \mathrm{AVG\_ENERGY}$ is roughly the work lost to flow friction per cycle. \codeword{pV.png} plots the last \codeword{--pvCycles} cycles (an unfinished last cycle is marked): the p--V loops with their work, p and $V$ against the crank angle, and the loop shape as $p/(C V^{-n}) - 1$ with $C$, $n$ fitted to each cycle. The pressure is the area average on the faces of the body whose motion table follows the gas volume (the piston crown, so the loop area is the piston work), or $\bar{p}$ when no such body is found.

\paragraph{Mass drift.}
A closed volume must keep its mass:
\begin{equation}
\mathrm{MASS\_DRIFT} = \frac{m(t_{end})}{m(0)} - 1 .
\end{equation}
\codeword{cycleConvergence.png} shows $m(t_k)/m(0) - 1$ at the end of every cycle. A drift comes from the mesh motion (space conservation) or a loose pressure tolerance, not from the physics. Values around $10^{-6}$ are fine. A drift of $10^{-3}$ shifts the mean pressure by the same fraction.

\subsubsection{Heat, Through-Flow and Porous Zones}
\paragraph{Heat per BC group.}
\codeword{wallHeatFlux} writes, per patch and time step, the integral of the heat flux over the patch, $\dot{Q}_{patch}$ (W, positive into the gas). For a group $g$ with the patches in \codeword{BC_GROUPS_RESOLVED}:
\begin{equation}
\mathrm{HEAT}_g = \frac{1}{t_1 - t_0}\int_{t_0}^{t_1}\sum_{patch \in g} \dot{Q}_{patch}\, dt \quad \textrm{(W)},
\end{equation}
and in \codeword{cycleConvergence.png} the heat per cycle $Q_{g,k} = \int_{t_{k-1}}^{t_k} \sum \dot{Q}\,dt$ (J). Energy balance of a closed engine at a periodic state, per cycle:
\begin{equation}
\sum_g Q_g \approx \mathrm{AVG\_ENERGY}
\end{equation}
(heat in = work out), and the mean gas temperature per cycle stops changing.

\paragraph{Inlets and outlets.}
\codeword{flow_<patch>} writes $\sum_f \phi_f$, with $\phi$ the face flux out of the domain (kg/s compressible, m$^3$/s incompressible). For a group:
\begin{equation}
\dot{m}_g = -\, s_\phi \sum_{patch \in g} \overline{\textstyle\sum_f \phi_f} ,
\end{equation}
positive into the domain, with $s_\phi = 1$ for compressible and $s_\phi = \rho_{ref}$ (\codeword{RHO_REF}, the gas density at the charge state) for incompressible runs. The pressure and temperature of a group are area weighted over its patches, with the patch areas $A$ from the header of the \codeword{avg_<patch>} files:
\begin{equation}
P_g = \frac{\sum A\,(p_0 + s_p\,\overline{p})}{\sum A}, \qquad T_g = \frac{\sum A\,\overline{T}}{\sum A} ,
\end{equation}
$s_p = 1$, $p_0 = 0$ for compressible, $s_p = \rho_{ref}$, $p_0$ = \codeword{CHARGE_PRESSURE} for incompressible (kinematic p to Pa). Then
\begin{equation}
\mathrm{MDOT\_IMBALANCE} = \frac{\sum_g \dot{m}_g}{\sum_g \max(\dot{m}_g, 0)}, \qquad
\mathrm{P\_DROP} = \mathrm{mean}_{inlets}\, P_g - \mathrm{mean}_{outlets}\, P_g .
\end{equation}
The imbalance is zero for a converged steady run, and small (the change of gas mass inside) for a transient one. Without \codeword{RHO_REF} an incompressible run reports \codeword{VOLUME_FLOW_<group>} in m$^3$/s and the kinematic pressure.

\paragraph{Porous zones.}
\codeword{porHeatToGas_<zone>} writes $\int_{zone} Q_s\,dV$ with $Q_s = h a\,(T_s - T)$ (Section~\ref{sec:regenModel}); its time average is \codeword{HEAT_<zone>} (W, from the matrix into the gas). \codeword{por_<zone>} writes the volume averages over the zone; their time averages are \codeword{P_MEAN_<zone>}, \codeword{T_MEAN_<zone>} and \codeword{TS_MEAN_<zone>} (matrix temperature).

\paragraph{Steady runs.}
The time is the iteration number. The averages above are taken over the last \codeword{avgIters} = 50 iterations (fewer if the run is shorter): $t_0 = \max(t_{first}, t_{end} - 50)$, $t_1 = t_{end}$. There are no cycles, work or p--V values. \codeword{flowConvergence.png} shows the initial residuals and the flow, pressure and temperature of each inlet/outlet group per iteration.

\subsubsection{Mesh Motion Numbers (Setup Report)}
These are computed by the setup before meshing, from the STL PIDs and the motion CSV.

\paragraph{Zone positions.}
Each motion zone is a flat face at height $z_{ref}$ (as in the STL). A zone of body $b$ is at
\begin{equation}
z(\theta) = z_{ref} + x_b(\theta) - x_b(\theta_{ref}) + \delta z_b ,
\end{equation}
with $x_b(\theta)$ the CSV column (periodic linear interpolation, any datum), $\theta_{ref}$ = \codeword{REF_ANGLE} (the angle the STL was exported at) and $\delta z_b$ the clearance offset (\codeword{CLEARANCE_MODE = shift}, else 0). Anchor zones stay at $z_{ref}$. The gap between two neighbouring zones (sorted by their mean height over the cycle) is $g(\theta) = z_{upper}(\theta) - z_{lower}(\theta)$, checked on a grid of 1° or finer. Any $g \le 0$ is an error (the zones cross).

\paragraph{Motion tables.}
For a run starting at $\theta_s$ (\codeword{START_ANGLE}, the meshing angle):
\begin{equation}
\theta(t) = \theta_s + 360 f t, \qquad \Delta z_b(t) = x_b(\theta(t)) - x_b(\theta_s) ,
\end{equation}
so every table starts at zero. Crank tables store $\Delta z_b(\theta)$ over 0 to 360° instead.

\paragraph{Clearance.}
A body with \codeword{TDC_CLEARANCE} $\delta$ must keep $g \ge \delta$ to each neighbouring anchor. The shortfall at the low end of the profile is $r = \delta - \min g$ (and the same at the high end, $l$).
\begin{itemize}
    \item \codeword{shift}: the whole path moves by $r$, $\delta z_b = r$. Added dead volume $r A$, with $A$ the projected zone area.
    \item \codeword{scale}: the stroke is shrunk to fit, keeping the shape exactly. With $x_{lo}$, $x_{hi}$ the profile ends,
    \begin{equation}
    x' = x_{lo} + r + (x - x_{lo})\,\frac{(x_{hi} - l) - (x_{lo} + r)}{x_{hi} - x_{lo}} .
    \end{equation}
    \item \codeword{soft} (default): only the part of the stroke near the closing end changes. With $s = x - x_{lo}$ and the stroke $S = x_{hi} - x_{lo}$,
    \begin{equation}
    x' = x_{lo} + \left(s^4 + r^4\right)^{1/4} - \epsilon\left(\frac{s}{S}\right)^4, \qquad \epsilon = \left(S^4 + r^4\right)^{1/4} - S .
    \end{equation}
    At $s = 0$ the end is raised by exactly $r$. Away from the end the change dies out as $r^4/(4s^3)$, and the small correction $\epsilon$ keeps the other end of the stroke exactly where it was. The profile stays monotonic and dwells stay flat. The peak acceleration near the adjusted end is about twice the original. The high end uses the mirror image. The report gives the largest change and over how many degrees the profile moved by more than 0.02 mm.
\end{itemize}
The added dead volume is $(r + l) A$ for the soft and scale modes too.

\paragraph{Squash and stretch.}
Between two zones the cells only scale along z, by $g(\theta)/g_0$ with $g_0 = g(\theta_s)$, the gap at the meshing angle. The worst ratio over the cycle and all gaps that change:
\begin{equation}
R = \max_{gaps}\, \max_\theta\, \max\!\left(\frac{g(\theta)}{g_0}, \frac{g_0}{g(\theta)}\right) .
\end{equation}
\codeword{START_ANGLE = auto} picks the grid angle with the smallest $R$, and \codeword{MAX_DEFORMATION} is the upper limit for $R$. For one gap that runs from $g_{min}$ to $g_{max}$, $R$ is smallest when the mesh is built at the geometric mean, where squash and stretch are equal:
\begin{equation}
g_0 = \sqrt{g_{min}\, g_{max}}, \qquad R = \sqrt{\frac{g_{max}}{g_{min}}} .
\end{equation}
For a piston with clearance $\delta$ and stroke $S$: $g_{min} = \delta$, $g_{max} = \delta + S$.

\paragraph{Body kinematics.}
On the angle grid, with $\omega_{deg} = 360 f$:
\begin{equation}
v = \frac{dz}{d\theta}\,\omega_{deg}, \qquad a = \frac{dv}{d\theta}\,\omega_{deg}
\end{equation}
(central differences). The report gives the stroke, $\max|v|$ and $\max|a|$, and warns when the velocity jumps by more than 20\,\% of $\max|v|$ between two samples (a sharp corner in the profile, which gives pressure spikes).

\paragraph{Gap refinement.}
For a side gap of width $w$ (smallest distance between the \codeword{GAP_PIDS} faces), base cell size $\Delta_0$ and \codeword{GAP_CELLS} $n$:
\begin{equation}
L = \left\lceil \log_2 \frac{n\,\Delta_0}{w} \right\rceil, \qquad \Delta = \frac{\Delta_0}{2^L} ,
\end{equation}
refined within a distance $1.25\,w$ of the gap faces.

\paragraph{Motion check frames.}
The dry run writes \codeword{MOTION_FRAMES} $n_f$ meshes per cycle (a multiple of 8) and runs $n_f\lceil 72/n_f\rceil$ steps per cycle, $\Delta t = 1/(f \cdot \textrm{steps})$. \codeword{checkMesh} runs at $t = k/(8f)$, $k = 1 \ldots 8$, every 45°.

\subsubsection{Worked Example}
A synthetic case that can be checked by hand. Piston with crown area $A = 0.01$\,m$^2$, gas below it, stroke $S = 0.05$\,m, $f = 10$\,Hz, position $z = \frac{S}{2}(1 - \cos\theta)$, so $V = V_c + A z$. The crown pressure peaks $\varphi$ after the smallest volume ($\theta = 0$), so it is high while the gas expands:
\begin{equation}
p = p_m + p_a \cos(\theta - \varphi), \qquad p_a = 20\,\textrm{kPa}, \quad \varphi = 60^\circ .
\end{equation}
\begin{itemize}
    \item Force on the crown $F_z = p A$, velocity $v = \frac{S}{2}\omega\sin\theta$. Work per cycle:
    \begin{equation}
    W = \int_0^{2\pi} p A \frac{S}{2}\sin\theta\, d\theta = \frac{\pi}{2} A S p_a \sin\varphi = \frac{\pi}{2}\cdot 0.01 \cdot 0.05 \cdot 20000 \cdot 0.866 = 13.6\,\textrm{J}.
    \end{equation}
    $p_m$ drops out, as said above. \codeword{AVG_ENERGY} = 13.6\,J, \codeword{AVG_POWER} = $13.6 \cdot 10 = 136$\,W. With the peak before the smallest volume ($\varphi < 0$, high pressure while the gas is compressed) the work is negative.
    \item p--V: with $\bar{p} = p$ the loop integral $\oint p\,dV = A\oint p\,dz$ is the same integral, so \codeword{PV_WORK} = 13.6\,J. In a real engine the face pressure and $\bar{p}$ differ and so do the two values.
    \item \codeword{FORCE_Z_AMP} $= p_a A = 200$\,N, \codeword{P_<crown>_AMP} = 20\,kPa.
    \item Crank drive: $T = F\,dz/d\theta = p A \frac{S}{2}\sin\theta$, $\oint T\,d\theta = 13.6$\,J, \codeword{TORQUE_GAS_MEAN} $= 13.6/(2\pi) = 2.165$\,N\,m.
    \item Generator $k = 1$\,V\,s/rad, $R_{load} = 10\,\Omega$, no friction. At 10\,Hz ($\omega = 62.8$\,rad/s): $I_e = 6.28$\,A, load torque 6.28\,N\,m, $P_{elec} = 395$\,W, more than the 136\,W the gas gives, so the crank slows down. If the loop stayed the same, the speed settles where the mean torques match, $2.165 = k^2\omega/R_{load}$: $\omega = 21.7$\,rad/s (3.45\,Hz), $P_{elec} = 2.165 \cdot 21.7 = 47$\,W = $W \cdot f$ at 3.45\,Hz.
    \item Gas volume: \codeword{gasMass} $m = 7.897\cdot10^{-3}$\,kg and \codeword{gasAvg} $\bar{\rho} = 1.1640$\,kg/m$^3$ give $V = 6.785\cdot10^{-3}$\,m$^3$ (6.8\,L).
    \item Motion: clearance $\delta = 2$\,mm, stroke 152\,mm: $g_0 = \sqrt{2 \cdot 154} = 17.5$\,mm, $R = \sqrt{154/2} = 8.8$. Soft clearance $r = 2$\,mm: at 10\,mm from the end the profile moves by $(10^4 + 2^4)^{1/4} - 10 = 0.004$\,mm, and $\epsilon = 1.1\cdot10^{-9}$\,m.
    \item Gap: $w = 2$\,mm, $\Delta_0 = 10.67$\,mm, $n = 6$: $\log_2(6 \cdot 10.67/2) = 5.0$, so $L = 5$, $\Delta = 0.33$\,mm.
\end{itemize}

\subsection{Demo Cases}
\codeword{internalDemoCase/CASES} has two cases (geometry in \codeword{internalDemoCase/02_reference/MSH}):
\begin{itemize}
    \item \codeword{staticDemo} - closed box, hot floor 600 K, cold ceiling 300 K, adiabatic sides, no motion.
    \item \codeword{motionDemo} - displacer with a dwell profile, 2 mm side gap and one hole, plus a power piston crown (\codeword{FROM = WALL-engine}), motion from \codeword{engineMotion.csv}.
\end{itemize}

\subsection{Limitations}
\begin{itemize}
    \item Motion along z only, bodies in line (no alpha engines), \codeword{displacementInterpolation} only.
    \item Sample surfaces (\codeword{SMP}) are not available yet. Regenerators and heater/cooler cores are porous zones (Section \ref{sec:porous}); a linear initial matrix temperature along the holes is not available yet.
    \item No \codeword{--continue} option yet. With the crank drive, raising \codeword{N_CYCLES} and rewriting the controlDict (\codeword{-d}) is enough, since the crank tables cover one cycle and repeat. With prescribed motion the motion tables only cover \codeword{N_CYCLES}, so the full setup has to be rerun as well.
    \item No p--V or cycle-convergence plots yet, only the summary and the flow images.
\end{itemize}
